\documentclass[fs13,seriesSau]{korfarm-base}
\usepackage{korfarm-saussure}

% ============================================================
% passage-note.tex (v11)
% 지문 박스 + 첫 페이지 우측 메모란 (동적 높이)
% 정책: PAGE_BREAK_POLICY.md §3, §4 참조
%
% 변경 (v11):
%  · 지문 박스 폭 확대: enlarge right by=-3.7cm  (메모 3.4cm + gap 3mm)
%  · 메모란 동적 높이: frame.north east ~ frame.south east 사이 자동 채움
%  · 메모 줄선: clip 영역으로 박스 내부만 채움 (박스 높이 모르고도 작동)
%  · 문단 들여쓰기 활성화: tcolorbox 내부 \parindent=1em, \noindent 제거
%
% 사용:
%   \kfPassageNote{제목}{지문 본문}      → 메모란 포함
%   \kfPassageNoteNT{지문 본문}          → 메모란 포함, 제목 없음
%   \kfPassageFull{제목}{본문}           → 메모란 없음 (활동 내 보기)
% ============================================================

\makeatletter

% 메모란 규격 (cm 고정)
\newlength{\kf@memoW}
\setlength{\kf@memoW}{3.4cm}
\newlength{\kf@memoGap}
\setlength{\kf@memoGap}{3mm}

% --- 메인: 제목 있는 버전 ---
\newcommand{\kf@passagenote@with}[2]{%
  \par\ifkfPassageNewPage\ifdim\pagetotal>0.4\textheight\clearpage\fi\fi\smallskip\needspace{6\baselineskip}%
  \begin{tcolorbox}[
    enhanced, breakable,
    width=\dimexpr\linewidth-4cm\relax,
    colback=kfPassageBg, colframe=kfPrimary!70,
    boxrule=0.6pt, arc=3pt,
    left=\kfBoxPad, right=\kfBoxPad, top=\dimexpr\kfBoxPad+8pt\relax, bottom=\kfBoxPad,
    title={\small\bfseries\color{kfPrimary} #1},
    coltitle=kfPrimary,
    colbacktitle=kfPrimaryLight!60,
    attach boxed title to top left={yshift=-2.4mm, xshift=8pt},
    boxed title style={colframe=kfPrimary!50, boxrule=0.5pt, arc=2pt},
    parbox=false,
    overlay unbroken and first={%
      \begin{scope}
        \draw[kfRule, line width=0.4pt, fill=kfNoteBg, rounded corners=2pt]
          ([xshift=\kf@memoGap]frame.north east) rectangle
          ([xshift=\kf@memoGap+\kf@memoW]frame.south east);
        \node[anchor=north east, font=\footnotesize\bfseries\color{kfMute}, inner sep=3pt]
          at ([xshift=\kf@memoGap+\kf@memoW, yshift=-2pt]frame.north east) {메모};
        \node[anchor=north east, fill=kfPrimary!15, text=kfPrimary,
              font=\scriptsize\bfseries, rounded corners=2pt,
              inner xsep=4pt, inner ysep=2pt]
          at ([xshift=\kf@memoGap+\kf@memoW-3pt, yshift=-19pt]frame.north east) {\kf@memocat};
        \begin{scope}
          \path[clip]
            ([xshift=\kf@memoGap]frame.north east) rectangle
            ([xshift=\kf@memoGap+\kf@memoW]frame.south east);
          \foreach \i in {1,...,40}{%
            \draw[kfRule, line width=0.25pt]
              ([xshift=\kf@memoGap+2mm, yshift=-7mm - \i*5mm]frame.north east) --
              ([xshift=\kf@memoGap+\kf@memoW-2mm, yshift=-7mm - \i*5mm]frame.north east);%
          }
        \end{scope}
      \end{scope}
    }
  ]
    \setlength{\parindent}{1em}%
    \hspace*{1em}\ignorespaces#2%
  \end{tcolorbox}\par\smallskip
}

% --- 제목 없는 버전 ---
\newcommand{\kf@passagenote@notitle}[1]{%
  \par\ifkfPassageNewPage\ifdim\pagetotal>0.4\textheight\clearpage\fi\fi\smallskip\needspace{6\baselineskip}%
  \begin{tcolorbox}[
    enhanced, breakable,
    width=\dimexpr\linewidth-4cm\relax,
    colback=kfPassageBg, colframe=kfPrimary!70,
    boxrule=0.6pt, arc=3pt,
    left=\kfBoxPad, right=\kfBoxPad, top=\dimexpr\kfBoxPad+8pt\relax, bottom=\kfBoxPad,
    parbox=false,
    overlay unbroken and first={%
      \begin{scope}
        \draw[kfRule, line width=0.4pt, fill=kfNoteBg, rounded corners=2pt]
          ([xshift=\kf@memoGap]frame.north east) rectangle
          ([xshift=\kf@memoGap+\kf@memoW]frame.south east);
        \node[anchor=north east, font=\footnotesize\bfseries\color{kfMute}, inner sep=3pt]
          at ([xshift=\kf@memoGap+\kf@memoW, yshift=-2pt]frame.north east) {메모};
        \node[anchor=north east, fill=kfPrimary!15, text=kfPrimary,
              font=\scriptsize\bfseries, rounded corners=2pt,
              inner xsep=4pt, inner ysep=2pt]
          at ([xshift=\kf@memoGap+\kf@memoW-3pt, yshift=-19pt]frame.north east) {\kf@memocat};
        \begin{scope}
          \path[clip]
            ([xshift=\kf@memoGap]frame.north east) rectangle
            ([xshift=\kf@memoGap+\kf@memoW]frame.south east);
          \foreach \i in {1,...,40}{%
            \draw[kfRule, line width=0.25pt]
              ([xshift=\kf@memoGap+2mm, yshift=-7mm - \i*5mm]frame.north east) --
              ([xshift=\kf@memoGap+\kf@memoW-2mm, yshift=-7mm - \i*5mm]frame.north east);%
          }
        \end{scope}
      \end{scope}
    }
  ]
    \setlength{\parindent}{1em}%
    \hspace*{1em}\ignorespaces#1%
  \end{tcolorbox}\par\smallskip
}

\newcommand{\kfPassageNote}[2]{%
  \def\kf@tmptitle{#1}%
  \ifx\kf@tmptitle\@empty
    \kf@passagenote@notitle{#2}%
  \else
    \kf@passagenote@with{#1}{#2}%
  \fi
}
\newcommand{\kfPassageNoteNT}[1]{\kf@passagenote@notitle{#1}}
\makeatother

% ============================================================
% 풀폭 지문 박스 (메모란 없음) — 활동 내 보기 박스
% PAGE_BREAK_POLICY.md §3
% ============================================================
\makeatletter
\newcommand{\kf@passagefull@with}[2]{%
  \par\ifkfPassageNewPage\ifdim\pagetotal>0.4\textheight\clearpage\fi\fi\smallskip\needspace{6\baselineskip}%
  \begin{tcolorbox}[
    enhanced, breakable,
    colback=kfPassageBg, colframe=kfPrimary!70,
    boxrule=0.6pt, arc=3pt,
    left=\kfBoxPad, right=\kfBoxPad, top=\dimexpr\kfBoxPad+8pt\relax, bottom=\kfBoxPad,
    title={\small\bfseries\color{kfPrimary} #1},
    coltitle=kfPrimary,
    colbacktitle=kfPrimaryLight!60,
    attach boxed title to top left={yshift=-2.4mm, xshift=8pt},
    boxed title style={colframe=kfPrimary!50, boxrule=0.5pt, arc=2pt},
    parbox=false
  ]
    \setlength{\parindent}{1em}%
    \hspace*{1em}\ignorespaces#2%
  \end{tcolorbox}\par\smallskip
}
\newcommand{\kf@passagefull@notitle}[1]{%
  \par\ifkfPassageNewPage\ifdim\pagetotal>0.4\textheight\clearpage\fi\fi\smallskip\needspace{6\baselineskip}%
  \begin{tcolorbox}[
    enhanced, breakable,
    colback=kfPassageBg, colframe=kfPrimary!70,
    boxrule=0.6pt, arc=3pt,
    left=\kfBoxPad, right=\kfBoxPad, top=\dimexpr\kfBoxPad+8pt\relax, bottom=\kfBoxPad,
    parbox=false
  ]
    \setlength{\parindent}{1em}%
    \hspace*{1em}\ignorespaces#1%
  \end{tcolorbox}\par\smallskip
}
\newcommand{\kfPassageFull}[2]{%
  \def\kf@tmptitle{#1}%
  \ifx\kf@tmptitle\@empty
    \kf@passagefull@notitle{#2}%
  \else
    \kf@passagefull@with{#1}{#2}%
  \fi
}
\newcommand{\kfPassageFullNT}[1]{\kf@passagefull@notitle{#1}}
\makeatother

% ============================================================
% 작은 문장 박스 (문장 독해용) — PAGE_BREAK_POLICY.md §2.5
% ============================================================
\newcommand{\kfSentenceBox}[2]{%
  \par\smallskip\needspace{3\baselineskip}%
  \noindent\begin{tcolorbox}[
    enhanced, breakable=false,
    colback=kfPassageBg, colframe=kfMute,
    boxrule=0.4pt, arc=2pt,
    left=8pt, right=6pt, top=4pt, bottom=4pt,
    title={\footnotesize\bfseries\color{kfPrimary} #1},
    coltitle=kfPrimary, colbacktitle=kfPaper,
    attach boxed title to top left={yshift=-1.6mm, xshift=4pt},
    boxed title style={colframe=kfMute, boxrule=0.3pt, arc=1pt}
  ]%
    \setstretch{1.3}#2%
  \end{tcolorbox}\par\smallskip%
}

% ============================================================
% 지문 본문 아래 안내/정리 박스 (v16 신규)
% 사용: \kfPassageHint{한 줄 정리}{본문 안내 텍스트}
% ============================================================
\newcommand{\kfPassageHint}[2]{%
  \par\smallskip\needspace{4\baselineskip}%
  \begin{tcolorbox}[
    enhanced, breakable=false,
    colback=kfPrimary!5, colframe=kfPrimary!40,
    boxrule=0.4pt, arc=2pt,
    left=8pt, right=6pt, top=4pt, bottom=4pt,
    title={\footnotesize\bfseries\color{kfPrimary} #1},
    coltitle=kfPrimary, colbacktitle=kfPaper,
    attach boxed title to top left={yshift=-1.5mm, xshift=4pt},
    boxed title style={colframe=kfPrimary!40, boxrule=0.3pt, arc=1pt}
  ]%
    \small\setstretch{1.3}#2%
  \end{tcolorbox}\par\smallskip%
}

% ============================================================
% 환경 버전 (호환성) — 메모란 포함
% ============================================================
\makeatletter
\newenvironment{kfPassageNoteEnv}[1]%
  {\par\needspace{6\baselineskip}%
   \def\kf@psrc{#1}%
   \begin{tcolorbox}[
     enhanced, breakable,
     width=\dimexpr\linewidth-4cm\relax,
     colback=kfPassageBg, colframe=kfPrimary!70,
     boxrule=0.6pt, arc=3pt,
     left=\kfBoxPad, right=\kfBoxPad, top=\dimexpr\kfBoxPad+8pt\relax, bottom=\kfBoxPad,
     title={\small\bfseries\color{kfPrimary} \kf@psrc},
     coltitle=kfPrimary, colbacktitle=kfPrimaryLight!60,
     attach boxed title to top left={yshift=-2.4mm, xshift=8pt},
     boxed title style={colframe=kfPrimary!50, boxrule=0.5pt, arc=2pt},
     parbox=false,
     overlay unbroken and first={%
       \begin{scope}
         \draw[kfRule, line width=0.4pt, fill=kfNoteBg, rounded corners=2pt]
           ([xshift=\kf@memoGap]frame.north east) rectangle
           ([xshift=\kf@memoGap+\kf@memoW]frame.south east);
         \node[anchor=north east, font=\footnotesize\bfseries\color{kfMute}, inner sep=3pt]
           at ([xshift=\kf@memoGap+\kf@memoW, yshift=-2pt]frame.north east) {메모};
         \begin{scope}
           \path[clip]
             ([xshift=\kf@memoGap]frame.north east) rectangle
             ([xshift=\kf@memoGap+\kf@memoW]frame.south east);
           \foreach \i in {1,...,40}{%
             \draw[kfRule, line width=0.25pt]
               ([xshift=\kf@memoGap+2mm, yshift=-7mm - \i*5mm]frame.north east) --
               ([xshift=\kf@memoGap+\kf@memoW-2mm, yshift=-7mm - \i*5mm]frame.north east);%
           }
         \end{scope}
       \end{scope}
     }
   ]%
   \setlength{\parindent}{1em}%
  }%
  {\end{tcolorbox}\par\smallskip}
\makeatother

% ============================================================
% condition-box.tex
% <조건> 박스 컴포넌트 (tcolorbox)
% 사용:
%   \begin{kfCondition}
%     \item 첫째 조건
%     \item 둘째 조건
%   \end{kfCondition}
% ============================================================

\newenvironment{kfCondition}%
  {\par\smallskip\needspace{4\baselineskip}%
   \begin{tcolorbox}[
     enhanced, breakable=false,
     colback=kfPrimaryLight!40, colframe=kfPrimary,
     boxrule=0.5pt, arc=2pt,
     left=\kfBoxPad, right=\kfBoxPad, top=\kfBoxPad, bottom=\kfBoxPad,
     title={\small\bfseries\color{kfPrimary} \,조건\,},
     coltitle=kfPrimary, colbacktitle=kfPaper,
     attach boxed title to top left={yshift=-2mm, xshift=6pt},
     boxed title style={colframe=kfPrimary, boxrule=0.4pt, arc=1pt}
   ]%
   \begin{enumerate}[leftmargin=16pt, itemsep=1pt, topsep=1pt,
                    label=\textbf{\arabic*.}]}%
  {\end{enumerate}\end{tcolorbox}\par\smallskip}

% ============================================================
% evidence-box.tex
% <보기> 박스 컴포넌트
% 사용:
%   \begin{kfEvidence}
%     보기 본문 ...
%   \end{kfEvidence}
% ============================================================

\newenvironment{kfEvidence}%
  {\par\smallskip\needspace{4\baselineskip}%
   \begin{tcolorbox}[
     enhanced, breakable=false,
     colback=kfPaper, colframe=kfMute,
     boxrule=0.5pt, arc=2pt,
     left=\kfBoxPad, right=\kfBoxPad, top=\kfBoxPad, bottom=\kfBoxPad,
     title={\small\bfseries\color{kfInk} \,보기\,},
     coltitle=kfInk, colbacktitle=kfPrimaryLight!50,
     attach boxed title to top left={yshift=-2mm, xshift=6pt},
     boxed title style={colframe=kfMute, boxrule=0.4pt, arc=1pt}
   ]%
   \leavevmode\mbox{}\ignorespaces%
  }%
  {\end{tcolorbox}\par\smallskip}

% ============================================================
% choice-list.tex
% 5선지 (① ② ③ ④ ⑤) 자동 들여쓰기 컴포넌트
% 사용:
%   \begin{kfChoices}
%     \kfChoice 첫째 선택지
%     \kfChoice 둘째 선택지
%     \kfChoice 셋째 선택지
%     \kfChoice 넷째 선택지
%     \kfChoice 다섯째 선택지
%   \end{kfChoices}
% ============================================================

\newcounter{kfChoiceCnt}
\newcommand{\kfChoiceMark}[1]{%
  \ifcase#1\relax%
  \or ①\or ②\or ③\or ④\or ⑤\or ⑥\or ⑦\or ⑧\or ⑨%
  \fi
}

% 환경: 자동 번호
\newenvironment{kfChoices}%
  {\setcounter{kfChoiceCnt}{0}%
   \par\smallskip%
   \begin{list}{}{%
     \setlength{\leftmargin}{20pt}%
     \setlength{\itemindent}{0pt}%
     \setlength{\labelwidth}{18pt}%
     \setlength{\labelsep}{6pt}%
     \setlength{\itemsep}{2pt}%
     \setlength{\topsep}{1pt}%
     \setlength{\parsep}{0pt}%
     \renewcommand{\makelabel}[1]{##1\hss}%
   }}%
  {\end{list}\par\smallskip}

\newcommand{\kfChoice}{%
  \stepcounter{kfChoiceCnt}%
  \item[\textbf{\kfChoiceMark{\value{kfChoiceCnt}}}]\ignorespaces%
}

% --- 한 줄 5선지 (짧은 선택지용) ---
\newcommand{\kfChoicesInline}[5]{%
  \par\smallskip%
  \noindent\textbf{①}~#1\quad\textbf{②}~#2\quad\textbf{③}~#3\quad\textbf{④}~#4\quad\textbf{⑤}~#5%
  \par\smallskip%
}

% ============================================================
% answer-note.tex
% 답안 작성란 (노트 스타일, 줄친 영역)
% 사용:
%   \kfAnswerLines{5}        % 5줄 답안란
%   \kfAnswerNote{서술형 답안란}{8} % 라벨+8줄
% ============================================================

% 단일 줄 (필요 시 인라인)
\newcommand{\kfAnswerLine}{%
  \par\noindent\textcolor{kfRule}{\rule{\linewidth}{0.4pt}}\par
}

% N줄 답안란 (간단)
\newcommand{\kfAnswerLines}[1]{%
  \par\smallskip%
  \begin{minipage}{\linewidth}%
  \foreach \i in {1,...,#1}{%
    \textcolor{kfRule}{\rule{\linewidth}{0.4pt}}\\[6pt]%
  }%
  \end{minipage}%
  \par\smallskip%
}

% 라벨이 있는 노트형 답안란
\newcommand{\kfAnswerNote}[2]{%
  \par\smallskip\needspace{#2\baselineskip}%
  \begin{tcolorbox}[
    enhanced, breakable=false,
    colback=kfPaper, colframe=kfRule,
    boxrule=0.4pt, arc=2pt,
    left=\kfBoxPad, right=\kfBoxPad, top=\kfBoxPad, bottom=\kfBoxPad,
    title={\footnotesize\bfseries\color{kfMute} #1},
    coltitle=kfMute, colbacktitle=kfPaper,
    attach boxed title to top left={yshift=-1.5mm, xshift=6pt},
    boxed title style={colframe=kfRule, boxrule=0.3pt, arc=1pt}
  ]%
  \foreach \i in {1,...,#2}{%
    \textcolor{kfRule}{\rule{\linewidth}{0.3pt}}\\[6pt]%
  }%
  \end{tcolorbox}\par\smallskip%
}

% 짧은 빈칸 (단답형)
\newcommand{\kfBlank}[1][3cm]{\underline{\hspace{#1}}}
% 넓은 빈칸 (서술 한 줄)
\newcommand{\kfBlankWide}[1][6cm]{\underline{\hspace{#1}}}
% 괄호 빈칸
\newcommand{\kfBlankParen}{(\,\hspace{1cm}\,)}
% O/X 빈칸
\newcommand{\kfBlankOX}{(\,\hspace{1.5cm}\,)}

% ============================================================
% question-block.tex
% 문제 블록 (번호 배지 + 발문 + 보기/조건 + 선택지 + 답안란)
% 모든 구성을 하나의 minipage로 묶어 단/페이지 단절 금지
%
% 사용:
%   \begin{kfQBlock}{1}{객관식}
%     \kfQStem{다음 글에 대한 설명으로 가장 적절한 것은?}
%     \begin{kfEvidence} ... \end{kfEvidence}    % 선택
%     \begin{kfCondition} ... \end{kfCondition}  % 선택
%     \begin{kfChoices}
%       \kfChoice ...
%     \end{kfChoices}
%     \kfAnswerLines{2}                            % 선택
%   \end{kfQBlock}
% ============================================================

% 무단절 minipage로 묶고, 발문/보기/조건/선택지/답안란을 자유롭게 배치
% 사용자 피드백 #4: [객관식]/[서술형] 등 텍스트 라벨 제거 — 번호 배지만 표시
% #2(유형 라벨)는 호환성 인자로만 받고 출력하지 않음.
\newenvironment{kfQBlock}[2]%
  {\par\smallskip\needspace{6\baselineskip}%
   \noindent\begin{minipage}{\linewidth}%
   \samepage%
   \par\noindent
   \kfQNum{#1}\hspace{4pt}%
   \begin{minipage}[t]{\dimexpr\linewidth-30pt\relax}%
  }%
  {\end{minipage}\end{minipage}\par\smallskip}

% 문제 발문
\newcommand{\kfQStem}[1]{%
  \par\noindent#1\par\vspace{2pt}%
}

% 짧은 소문항 라벨 (1), (2), (3)
\newcounter{kfSubQ}
\newcommand{\kfSubQReset}{\setcounter{kfSubQ}{0}}
\newcommand{\kfSubQ}{%
  \stepcounter{kfSubQ}%
  \par\noindent\textbf{(\arabic{kfSubQ})}~\ignorespaces%
}

% ============================================================
% activity-table.tex
% 활동: 정리표 / 내용 확인표 (마크다운 표 → tabularx 변환)
% 사용:
%   % 일반 tabular (직접 컬럼 명시)
%   \begin{kfActTable}{|l|X|X|}
%     \kfTHead{항목} & \kfTHead{설명} & \kfTHead{학생 작성} \\\hline
%     ① & ... & \kfTBlank \\\hline
%   \end{kfActTable}
%
%   % adjustbox 자동 축소 (정해진 폭으로 강제)
%   \kfActTableWrap{
%     ... (\begin{tabular}...\end{tabular} 코드) ...
%   }
% ============================================================

% 사용 시 본문 마지막 행은 \\\hline 으로 끝나야 함.
% v17.1: 흰색 mask로 Canva 배경 본문 침범 차단
\newenvironment{kfActTable}[1]%
  {\par\smallskip\needspace{4\baselineskip}%
   \renewcommand{\arraystretch}{1.3}%
   \arrayrulecolor{kfRule}%
   \begin{tcolorbox}[blanker, colback=white, left=2pt, right=2pt, top=2pt, bottom=2pt]%
   \noindent\begin{adjustbox}{max width=\linewidth}%
   \begin{tabular}{#1}%
   \hline}%
  {\end{tabular}%
   \end{adjustbox}%
   \end{tcolorbox}%
   \arrayrulecolor{black}%
   \par\smallskip}

% 헤더 행 헬퍼
\newcommand{\kfTHead}[1]{\textbf{\color{kfPrimary} #1}}
\newcommand{\kfTHeadRow}[1]{\rowcolor{kfPrimaryLight!50}\textbf{\color{kfPrimary} #1}}

% 빈 셀 (학생 작성용)
\newcommand{\kfTBlank}{\rule[-1.6em]{0pt}{3.4em}}
\newcommand{\kfTBlankSm}{\rule[-1.0em]{0pt}{2.4em}}

% adjustbox 강제 축소 래퍼 — Python에서 변환된 raw tabular 블록을 감쌀 때 사용
\newcommand{\kfActTableWrap}[1]{%
  \par\smallskip\needspace{4\baselineskip}%
  \begin{tcolorbox}[blanker, colback=white, left=2pt, right=2pt, top=2pt, bottom=2pt]%
  \noindent\begin{adjustbox}{max width=\linewidth, center}%
  #1%
  \end{adjustbox}%
  \end{tcolorbox}\par\smallskip%
}

% ============================================================
% activity-vocab.tex
% 어휘 학습 표 (3열: 번호 - 뜻 - 빈칸)
% 사용:
%   \begin{kfVocab}
%     \kfVocabItem{1}{눈으로 보고 마음으로 깨달음}
%     \kfVocabItem{2}{사물의 이치를 따져 헤아림}
%   \end{kfVocab}
% ============================================================

% adjustbox로 폭 강제 + 일반 tabular + p{} 열 사용
\newenvironment{kfVocab}%
  {\par\smallskip\needspace{4\baselineskip}%
   \renewcommand{\arraystretch}{1.35}%
   \arrayrulecolor{kfRule}%
   \noindent\begin{adjustbox}{max width=\linewidth, center}%
   \begin{tabular}{|>{\centering\arraybackslash}m{1.0cm}|m{8.5cm}|>{\centering\arraybackslash}m{4.0cm}|}%
   \hline
   \rowcolor{kfPrimaryLight!50}%
   \multicolumn{1}{|c|}{\textbf{\color{kfPrimary}번호}} &
   \multicolumn{1}{c|}{\textbf{\color{kfPrimary}뜻풀이}} &
   \multicolumn{1}{c|}{\textbf{\color{kfPrimary}어휘 (정답 작성)}} \\\hline
   \ignorespaces}%
  {\end{tabular}%
   \end{adjustbox}%
   \arrayrulecolor{black}%
   \par\smallskip}

\newcommand{\kfVocabItem}[2]{%
  \textbf{#1} & #2 & \rule[-1.0em]{0pt}{2.4em}\\\hline%
}

% ============================================================
% activity-blank.tex
% 빈칸 채우기 활동
% 사용:
%   \begin{kfBlankList}
%     \kfBlankItem 첫 번째 문장에서 (\kfBlank)은(는) ...
%     \kfBlankItem 둘째 문장에서 ...
%   \end{kfBlankList}
% ============================================================

\newenvironment{kfBlankList}%
  {\par\smallskip%
   \begin{enumerate}[leftmargin=20pt, itemsep=4pt, topsep=2pt,
                    label=\textbf{\arabic*.}, labelsep=6pt]}%
  {\end{enumerate}\par\smallskip}

\newcommand{\kfBlankItem}{\item}

% 텍스트 안에 박스형 빈칸 (단어 1개 채우기)
\newcommand{\kfBlankBox}[1][2.4cm]{%
  \tikz[baseline=-0.5ex]{%
    \draw[kfRule, line width=0.4pt] (0,0) -- ++(#1,0);%
  }%
}

% ============================================================
% activity-ox.tex (v15)
% O/X 표기 활동 — 그래픽 디자인
% 사용:
%   \begin{kfOXList}
%     \kfOXItem 글쓴이는 자연을 예찬하고 있다.\kfOX
%     \kfOXItem 1연과 4연은 수미상관 구조이다.\kfOX
%   \end{kfOXList}
%
% \kfOX = 우측에 [O] [X] 두 박스 (학생이 하나 동그라미)
% ============================================================

\newenvironment{kfOXList}%
  {\par\smallskip%
   \begin{enumerate}[leftmargin=20pt, itemsep=5pt, topsep=2pt,
                    label=\textbf{\arabic*.}, labelsep=6pt]}%
  {\end{enumerate}\par\smallskip}

\newcommand{\kfOXItem}{\item}

% v16: O/X 그래픽 박스 키움 (18×14 → 24×20pt) + 영역색 강조
\newcommand{\kfOX}{%
  \hfill\nolinebreak
  \tikz[baseline=-0.9ex]{%
    \node[draw=kfPrimary, fill=kfPrimary!5, line width=0.6pt, rounded corners=3pt,
          minimum width=24pt, minimum height=20pt, inner sep=0pt,
          font=\normalsize\bfseries\color{kfPrimary}] (ob) {O};%
    \node[draw=kfMute, fill=kfMute!5, line width=0.6pt, rounded corners=3pt,
          minimum width=24pt, minimum height=20pt, inner sep=0pt,
          font=\normalsize\bfseries\color{kfMute},
          right=4pt of ob] (xb) {X};%
  }%
}

% ============================================================
% activity-connect.tex
% 좌우 선 긋기 활동 (2단 양식 + 중앙 여백)
% 사용:
%   \begin{kfConnect}
%     \kfPair{사르트르}{실존은 본질에 앞선다}
%     \kfPair{니체}{초인 사상}
%     \kfPair{하이데거}{현존재}
%   \end{kfConnect}
% ============================================================

\newenvironment{kfConnect}%
  {\par\smallskip%
   \renewcommand{\arraystretch}{1.6}%
   \begin{tabular}{@{}>{\raggedleft\arraybackslash}m{4.5cm}
                     @{\hspace{2.5cm}}
                     >{\raggedright\arraybackslash}m{6.5cm}@{}}}%
  {\end{tabular}\par\smallskip}

\newcommand{\kfPair}[2]{%
  \tikz[baseline]{\node[draw=kfRule, fill=kfPrimaryLight!30, rounded corners=2pt,
        inner sep=4pt, font=\small]{#1};} \tikz[baseline]{\node[circle, draw=kfMute, fill=kfPaper, inner sep=1.5pt]{};}
  &
  \tikz[baseline]{\node[circle, draw=kfMute, fill=kfPaper, inner sep=1.5pt]{};} \tikz[baseline]{\node[draw=kfRule, fill=kfNoteBg, rounded corners=2pt,
        inner sep=4pt, font=\small]{#2};} \\
}

% ============================================================
% activity-writing.tex
% 글쓰기 활동 (큰 노트 영역)
% 사용:
%   \kfWriting{독후감}{12}     % 라벨 + 12줄 큰 노트
%   \kfWritingPrompt{프롬프트 텍스트}
% ============================================================

\newcommand{\kfWritingPrompt}[1]{%
  \par\smallskip%
  \begin{tcolorbox}[
    enhanced, breakable=false,
    colback=kfPrimaryLight!30, colframe=kfPrimary,
    boxrule=0.4pt, arc=3pt,
    left=\kfBoxPad, right=\kfBoxPad, top=\kfBoxPad, bottom=\kfBoxPad,
  ]%
  \small\textbf{\color{kfPrimary}글쓰기 안내}\\[2pt]%
  #1
  \end{tcolorbox}\par\smallskip%
}

\newcommand{\kfWriting}[2]{%
  \par\smallskip\needspace{\numexpr#2+2\relax\baselineskip}%
  \begin{tcolorbox}[
    enhanced, breakable=false,
    colback=kfPaper, colframe=kfRule,
    boxrule=0.4pt, arc=2pt,
    left=\kfBoxPad, right=\kfBoxPad, top=\kfBoxPad, bottom=\kfBoxPad,
    title={\footnotesize\bfseries\color{kfMute} #1},
    coltitle=kfMute, colbacktitle=kfPaper,
    attach boxed title to top left={yshift=-1.5mm, xshift=6pt},
    boxed title style={colframe=kfRule, boxrule=0.3pt, arc=1pt}
  ]%
  \foreach \i in {1,...,#2}{%
    \textcolor{kfRule}{\rule{\linewidth}{0.3pt}}\\[7pt]%
  }%
  \end{tcolorbox}\par\smallskip%
}

% ============================================================
% activity-structure.tex
% 구조도 (트리 + 빈칸)
% 사용:
%   \begin{kfStructure}
%     \kfNode{0}{서론}{문제 제기}
%     \kfNode{1}{본론}{(\,\quad\,)}
%     \kfNode{1}{본론}{근거 2}
%     \kfNode{0}{결론}{주장 강화}
%   \end{kfStructure}
% ============================================================

\newenvironment{kfStructure}%
  {\par\smallskip\needspace{6\baselineskip}%
   \begin{tcolorbox}[
     enhanced, breakable=false,
     colback=kfPaper, colframe=kfRule,
     boxrule=0.4pt, arc=2pt,
     left=\kfBoxPad, right=\kfBoxPad, top=\kfBoxPad, bottom=\kfBoxPad,
     title={\footnotesize\bfseries\color{kfMute} 글의 구조},
     coltitle=kfMute, colbacktitle=kfPrimaryLight!40,
     attach boxed title to top left={yshift=-1.5mm, xshift=6pt},
     boxed title style={colframe=kfRule, boxrule=0.3pt, arc=1pt}
   ]%
   \begin{tabbing}
     \hspace{1.4cm}\=\hspace{1.4cm}\=\hspace{8cm}\kill}%
  {\end{tabbing}\end{tcolorbox}\par\smallskip}

% 노드: #1=들여쓰기 단계 (0~3), #2=라벨, #3=내용
\newcommand{\kfNode}[3]{%
  \ifcase#1\relax%
    \>\>\textbf{\color{kfPrimary} ▶ #2}\quad #3\\
  \or
    \>\>\quad\textbf{\color{kfMute} ◦ #2}\quad #3\\
  \or
    \>\>\quad\quad\textbf{\color{kfMute} - #2}\quad #3\\
  \or
    \>\>\quad\quad\quad\textbf{\color{kfMute} \textperiodcentered\ #2}\quad #3\\
  \fi
}

% 빈칸 노드
\newcommand{\kfNodeBlank}[2]{\kfNode{#1}{#2}{(\hspace{2.5cm})}}

% ============================================================
% activity-sentence-reading.tex
% 문장 독해 활동 — 문장 박스 1개 + 그에 묶인 문제 N개를 한 minipage로
% 사용:
%   \begin{kfSentenceUnit}{1}{"바람은 내 귀에 속삭이며..."}
%     \kfSentQ{(1)}{서술어 '속삭이며'의 주어는?}
%     \begin{kfChoices}
%       \kfChoice 바람은
%       ...
%     \end{kfChoices}
%   \end{kfSentenceUnit}
% ============================================================

% kfSentenceBox 매크로는 passage-note.tex에 정의됨

\newenvironment{kfSentenceUnit}[2]%
  {\par\smallskip\needspace{6\baselineskip}%
   \noindent\begin{minipage}{\linewidth}%
   \samepage%
   \kfSentenceBox{문장 #1}{#2}%
   \par\smallskip%
  }%
  {\end{minipage}\par\smallskip}

% 같은 문장에 묶인 소문항 (문장 안내 + 발문)
\newcommand{\kfSentQ}[2]{%
  \par\noindent\textbf{#1}~#2\par\smallskip%
}

% 헤더 없이 문장만 있는 묶음 (sentence_ref가 없는 일반 문항)
\newenvironment{kfSentencelessUnit}%
  {\par\smallskip\noindent\begin{minipage}{\linewidth}\samepage}%
  {\end{minipage}\par\smallskip}

% ============================================================
% chapter-cover.tex (v6)
% 챕터 진입 페이지 (표지) — 하단에 영역 목차 추가
% 사용:
%   \kfChapterCover{1}{근대성과 자아}{Chapter 1}{이번 챕터에서는 ...}
%   \begin{kfChapterAreaList}
%     \kfChapterAreaItem{어휘}{kfAreaVocab}{핵심 어휘 12개}
%     ...
%   \end{kfChapterAreaList}
%
%   영역 목차가 없을 때는 \kfChapterCoverEnd 호출
% ============================================================

\newcommand{\kfChapterCover}[4]{%
  \clearpage%
  \thispagestyle{kfBlank}%
  \kfMarkChapter{Chapter #1. #2}%
  \kfChapterCoverBg%
  % v18.5: 챕터 표지 우상단에 로고
  \begin{tikzpicture}[remember picture, overlay]
    \node[anchor=north east, inner sep=0pt] at ($(current page.north east)+(-18mm,-18mm)$) {\kfLogoMd};
  \end{tikzpicture}%
  \vspace*{20mm}%
  \noindent
  \begin{tikzpicture}[remember picture, overlay]
    % 좌측 액센트 바
    \fill[kfPrimary] (current page.west) ++(12mm,25mm) rectangle ++(5mm, -50mm);
    % 챕터 번호 배지 (이중 원, 약간 작게)
    \node[anchor=north west, inner sep=0pt] at ($(current page.north west)+(24mm,-45mm)$) {%
      \begin{tikzpicture}
        \draw[kfPrimary, line width=1pt] (0,0) circle (10mm);
        \draw[kfPrimary, line width=0.4pt] (0,0) circle (8mm);
        \node[font=\normalsize\bfseries\color{kfPrimary}] at (0,2mm) {CHAPTER};
        \node[font=\LARGE\bfseries\color{kfPrimary}] at (0,-3mm) {#1};
      \end{tikzpicture}%
    };
  \end{tikzpicture}%
  \vspace{42mm}%
  \par\noindent\hspace{32mm}%
  \begin{minipage}{0.65\linewidth}%
    {\footnotesize\color{kfMute} #3}\\[4pt]
    {\LARGE\bfseries\color{kfPrimary} #2}\\[8pt]
    {\color{kfRule}\rule{50mm}{0.5pt}}\\[6pt]
    {\small\color{kfInk} #4}
  \end{minipage}%
  \par\vspace{14mm}%
}

% v6 / v18.5 / v18.6: 챕터 표지의 영역 목차 — 흰색 반투명 박스
% v18.6: 배경 가로선/도형과 안 겹치도록 TikZ overlay로 우하단 절대 위치
\makeatletter
% 영역 항목들을 모아 둘 박스
\newsavebox{\kf@arealistbox}
\newenvironment{kfChapterAreaList}{%
  \begin{lrbox}{\kf@arealistbox}%
  \begin{minipage}{0.62\linewidth}%
    \begin{tcolorbox}[%
      enhanced, colback=white, colframe=white,
      opacityback=0.92, opacityframe=0,
      boxrule=0pt, arc=4pt,
      width=\linewidth,
      top=8pt, bottom=8pt, left=10pt, right=10pt,
    ]%
      {\footnotesize\bfseries\color{kfMute} 이 챕터의 영역}\par\vspace{4pt}%
      {\color{kfRule}\hrule height 0.3pt}\par\vspace{4pt}%
      \begin{itemize}[leftmargin=14pt, itemsep=2pt, topsep=0pt, label={\color{kfPrimary!70}\textbullet}]%
}{%
      \end{itemize}%
    \end{tcolorbox}%
  \end{minipage}%
  \end{lrbox}%
  % v18.7: 배경 상단 가로선 ~ 하단 가로선 사이의 중간 영역에 배치
  % 가로선이 내용을 가리지 않도록 박스 중심이 두 선 사이 중점에 위치
  \begin{tikzpicture}[remember picture, overlay]%
    \node[anchor=center, inner sep=0pt]
      at ($(current page.south)+(0mm,90mm)$)
      {\usebox{\kf@arealistbox}};%
  \end{tikzpicture}%
  \vfill%
  \noindent\hfill\kfSeriesBadge\hspace{12mm}%
  \clearpage%
}
\makeatother

% 영역 항목: \kfChapterAreaItem{영역명}{kfArea색}{부제}
\makeatletter
\newcommand{\kfChapterAreaItem}[3]{%
  \item {\small\bfseries\color{#2} #1}%
  \def\kf@cci@sub{#3}%
  \ifx\kf@cci@sub\@empty\else
    {\small\color{kfMute}\, \textemdash\, #3}%
  \fi
}
\makeatother

% 영역 목록 없이 cover만 (legacy)
\newcommand{\kfChapterCoverEnd}{%
  \vfill%
  \noindent\hfill\kfSeriesBadge\hspace{12mm}%
  \clearpage%
}

% ============================================================
% area-header.tex (v16)
% 영역 시작 표제 — 시리즈별 차별화 라벨 + 큰 영역명 + 부제
% PAGE_BREAK_POLICY.md §1.4
%
% v16 (디자인 고도화 Phase 1):
%   - 시리즈별 라벨 박스 추가:
%     소쉬르 = AREA START (영역색 채움 박스)
%     프레게 = SECTION OPEN (영역색 테두리)
%     러셀   = RUSSELL (영역색 라이트 박스)
%     비트   = WITTGENSTEIN (회색 라이트 박스)
%   - 라벨 위치: 영역명 위 좌측
% ============================================================

\makeatletter

% 시리즈 라벨 텍스트
\newcommand{\kf@serieslabel}{%
  \ifcase\kf@series\relax
    AREA START%
  \or
    AREA START%
  \or
    SECTION OPEN%
  \or
    RUSSELL%
  \or
    WITTGENSTEIN%
  \fi
}

% 시리즈 라벨 박스 스타일 (#1 = 영역색)
\newcommand{\kf@labelbox}[1]{%
  \ifcase\kf@series\relax
    % 0/1 = 소쉬르: 영역색 채움
    \tikz[baseline=-2pt]{\node[fill=#1, text=white,
      rounded corners=3pt, inner xsep=6pt, inner ysep=2pt,
      font=\scriptsize\bfseries]{\kf@serieslabel};}%
  \or
    % 1 = 소쉬르: 영역색 채움
    \tikz[baseline=-2pt]{\node[fill=#1, text=white,
      rounded corners=3pt, inner xsep=6pt, inner ysep=2pt,
      font=\scriptsize\bfseries]{\kf@serieslabel};}%
  \or
    % 2 = 프레게: 영역색 테두리
    \tikz[baseline=-2pt]{\node[draw=#1, line width=0.6pt, text=#1,
      rounded corners=3pt, inner xsep=6pt, inner ysep=2pt,
      font=\scriptsize\bfseries]{\kf@serieslabel};}%
  \or
    % 3 = 러셀: 영역색 라이트 박스
    \tikz[baseline=-2pt]{\node[fill=#1!12, text=#1,
      rounded corners=2pt, inner xsep=5pt, inner ysep=1.5pt,
      font=\scriptsize\bfseries]{\kf@serieslabel};}%
  \or
    % 4 = 비트: 회색 미니 박스
    \tikz[baseline=-2pt]{\node[fill=kfMute!10, text=kfMute,
      rounded corners=2pt, inner xsep=5pt, inner ysep=1.5pt,
      font=\scriptsize\bfseries]{\kf@serieslabel};}%
  \fi
}

% v17: 시리즈+영역별 Canva 배경 자동 매핑
% \kf@hasbg=1로 set되면 \kfAreaStart에서 라벨 박스 출력 생략
\def\kf@areaVocab{어휘}\def\kf@areaGrammar{문법}\def\kf@areaConcept{개념}
\def\kf@areaLit{문학}\def\kf@areaNonlit{비문학}
\newif\ifkf@bgon
% 파일 있으면 배경 적용 + boolean true. 없으면 nothing.
\newcommand{\kfBgSet}[1]{%
  \IfFileExists{#1.pdf}{\kfPageBg{#1}\kf@bgontrue}{}%
}
\newcommand{\kf@trybg}[1]{%
  % v18.4: 영역 배경 PDF 비활성화 — 큰 도형이 본문 영역을 침범해
  % "영역 헤더 후 빈 페이지" 문제를 일으킴. 라벨 박스로만 영역 표시.
  \kf@bgonfalse%
}

% Note: 러셀(rus_*)/비트(wit_*) 배경 PDF 미존재 시 \kfPageBg가 에러 → \IfFileExists로 가드
% (현재는 파일이 모두 존재한다고 가정. 부분 제공 시 cls 수정 필요)

\newcommand{\kfAreaStart}[3]{%
  \clearpage%
  \thispagestyle{fancy}%
  \kfMarkArea{#1}%
  \kf@trybg{#1}%
  \vspace*{1mm}%
  % 배경 없으면 라텍스 라벨 박스 출력 (배경 있으면 일러스트가 라벨 역할)
  \ifkf@bgon
    \vspace*{4mm}%
  \else
    \noindent\kf@labelbox{#2}\par\vspace{4pt}%
  \fi
  % 영역명 + 부제
  \noindent
  \begin{tikzpicture}
    \node[anchor=west, inner sep=0pt] (bar) at (0,0) {\color{#2}\rule{3pt}{22pt}};
    \node[anchor=base west, inner sep=0pt, xshift=8pt, yshift=2pt] (lbl) at (bar.east |- 0,0) {%
      {\LARGE\bfseries\color{#2} #1}%
    };
    \def\kf@areasub{#3}%
    \ifx\kf@areasub\@empty\else
      \node[anchor=base west, inner sep=0pt, xshift=8pt, font=\normalsize\color{kfMute}]
        at (lbl.base east) {\,\textperiodcentered\, #3};%
    \fi
  \end{tikzpicture}%
  \par\vspace{2pt}
  {\color{#2!70}\hrule height 1pt}%
  \par\vspace{1pt}%
  {\color{kfRule}\hrule height 0.3pt}%
  \par\vspace{4pt}%
  \needspace{6\baselineskip}\nopagebreak[4]%
  \global\kf@areaJustOpenedtrue% v18.5: 첫 컨텐츠 needspace 무시
}
\makeatother

% 시리즈/영역 단축 매크로
\newcommand{\kfStartVocab}[1]{\kfAreaStart{어휘}{kfAreaVocab}{#1}}
\newcommand{\kfStartGrammar}[1]{\kfAreaStart{문법}{kfAreaGrammar}{#1}}
\newcommand{\kfStartConcept}[1]{\kfAreaStart{개념}{kfAreaConcept}{#1}}
\newcommand{\kfStartLit}[1]{\kfAreaStart{문학}{kfAreaLiterature}{#1}}
\newcommand{\kfStartNonlit}[1]{\kfAreaStart{비문학}{kfAreaNonlit}{#1}}
\newcommand{\kfStartWeekly}[1]{\kfAreaStart{실력 확인}{kfAreaWeekly}{#1}}

% ============================================================
% section-header.tex
% 섹션 시작 라벨 헤더 — [지문] / [활동] / [문제] 라벨
% 좌측 액센트 바 + 라벨 + 부제 (영역 액센트 색)
%
% 사용:
%   \kfSecHeader{kfAreaLiterature}{지문}{빼앗긴 들에도 봄은 오는가 / 이상화}
%   \kfSecHeader{kfAreaConcept}{활동}{내용 확인}
%   \kfSecHeader{kfAreaNonlit}{문제}{}
%
% 헤더 직후 페이지 분할 금지: \needspace{4\baselineskip} + \nopagebreak
% ============================================================

\providecommand{\kfSecHeader}[3]{%
  \par\needspace{10\baselineskip}\filbreak%
  \vspace{3pt}%
  \noindent
  \begin{tikzpicture}
    \node[anchor=west, inner sep=0pt] (bar) at (0,0) {\color{#1}\rule{3pt}{14pt}};
    \node[anchor=west, inner sep=0pt, xshift=6pt, font=\normalsize\bfseries\color{#1}]
       (lbl) at (bar.east) {[\,#2\,]};
    \node[anchor=west, inner sep=0pt, xshift=4pt, font=\small\color{kfMute}]
       at (lbl.east) {#3};
  \end{tikzpicture}%
  \par\vspace{1pt}%
  {\color{#1!50}\hrule height 0.4pt}%
  \par\vspace{2pt}\nopagebreak%
}

% 영역별 단축 매크로
\newcommand{\kfSecPassageVocab}[1]{\kfSecHeader{kfAreaVocab}{지문}{#1}}
\newcommand{\kfSecPassageGrammar}[1]{\kfSecHeader{kfAreaGrammar}{지문}{#1}}
\newcommand{\kfSecPassageConcept}[1]{\kfSecHeader{kfAreaConcept}{지문}{#1}}
\newcommand{\kfSecPassageLit}[1]{\kfSecHeader{kfAreaLiterature}{지문}{#1}}
\newcommand{\kfSecPassageNonlit}[1]{\kfSecHeader{kfAreaNonlit}{지문}{#1}}
\newcommand{\kfSecPassageWeekly}[1]{\kfSecHeader{kfAreaWeekly}{지문}{#1}}

\newcommand{\kfSecActVocab}[1]{\kfSecHeader{kfAreaVocab}{활동}{#1}}
\newcommand{\kfSecActGrammar}[1]{\kfSecHeader{kfAreaGrammar}{활동}{#1}}
\newcommand{\kfSecActConcept}[1]{\kfSecHeader{kfAreaConcept}{활동}{#1}}
\newcommand{\kfSecActLit}[1]{\kfSecHeader{kfAreaLiterature}{활동}{#1}}
\newcommand{\kfSecActNonlit}[1]{\kfSecHeader{kfAreaNonlit}{활동}{#1}}
\newcommand{\kfSecActWeekly}[1]{\kfSecHeader{kfAreaWeekly}{활동}{#1}}

\newcommand{\kfSecQVocab}[1]{\kfSecHeader{kfAreaVocab}{문제}{#1}}
\newcommand{\kfSecQGrammar}[1]{\kfSecHeader{kfAreaGrammar}{문제}{#1}}
\newcommand{\kfSecQConcept}[1]{\kfSecHeader{kfAreaConcept}{문제}{#1}}
\newcommand{\kfSecQLit}[1]{\kfSecHeader{kfAreaLiterature}{문제}{#1}}
\newcommand{\kfSecQNonlit}[1]{\kfSecHeader{kfAreaNonlit}{문제}{#1}}
\newcommand{\kfSecQWeekly}[1]{\kfSecHeader{kfAreaWeekly}{문제}{#1}}


\begin{document}

\thispagestyle{kfBlank}
\kfBookCoverBg
\vspace*{\kfBookCoverVspace}
\begin{center}
{\kfLogoLg}\\[14pt]
{\fontsize{72}{84}\selectfont\bfseries\kfBookCoverColor 소쉬르}\\[18pt]
{\fontsize{28}{32}\selectfont\kfBookCoverSubColor \textsc{Level} \textbf{2}}\\[22pt]
{\large\kfBookCoverSubColor 제 1 권 \quad\textbar\quad 챕터 1\textendash{} 4}
\end{center}
\vfill
\hfill\kfSeriesBadge\hspace{15mm}
\clearpage
\kfChapterCover{1}{소쉬르2 (20주차) (1)}{Chapter 1}{이번 챕터의 학습 내용을 시작합니다.}
\begin{kfChapterAreaList}
\kfChapterAreaItem{어휘}{kfAreaVocab}{나를 표현하는 낱말 익히기}
\kfChapterAreaItem{개념}{kfAreaConcept}{자기소개 방법}
\kfChapterAreaItem{문학}{kfAreaLiterature}{「나는 나야, 특별한 나」}
\kfChapterAreaItem{비문학}{kfAreaNonlit}{나를 알리는 자기소개 쓰기}
\kfChapterAreaItem{실력확인}{kfAreaWeekly}{}
\end{kfChapterAreaList}

\kfStartVocab{나를 표현하는 낱말 익히기}


\begin{kfVocab}
\kfVocabItem{1}{앞으로 이루고 싶은 희망}
\kfVocabItem{2}{함께 사는 친척들}
\kfVocabItem{3}{사람이 가지고 있는 마음의 특징}
\kfVocabItem{4}{좋은 점, 잘하는 점}
\kfVocabItem{5}{즐겨서 하는 일}
\kfVocabItem{6}{남보다 특별히 잘하는 것}
\end{kfVocab}


\kfStartConcept{자기소개 방법}


\kfPassageNoteNT{%
자기소개는 다른 사람에게 나를 알리는 중요한 일이에요. 자기소개를 잘 하려면 세 가지를 기억하세요.

먼저, 거짓말을 하지 않고 진실하게 말해야 해요.

그리고 "운동을 좋아해요"보다 "매주 토요일에 축구를 해요"처럼 구체적으로 말하는 것이 더 좋아요.

마지막으로, 나의 좋은 점을 자신 있게 긍정적으로 말해요.
\par\medskip\begin{itemize}[leftmargin=18pt, itemsep=2pt, topsep=2pt, label=\textbullet]
\item 진실하게: 거짓 없이 사실대로 말해요.
\item 구체적으로: '운동을 좋아해요' → '매주 토요일에 축구를 해요'
\item 긍정적으로: 좋은 점을 자신 있게 말해요.
\end{itemize}
}

\kfPassageNote{문장의 기본 구조: 주어 + 서술어}{%
모든 문장은 주어와 서술어로 이루어져 있어요. \\[4pt]
주어: '누가' 또는 '무엇이'에 해당하는 말 \\[4pt]
서술어: '어떠하다', '무엇이다', '어떻게 한다'에 해당하는 말 \\[4pt]
예시: 나는 (주어) 학생입니다. (서술어)
}


\kfActivityBg \par\kfMaybeNeedspace{24\baselineskip}\noindent{\kfTipMark\,\,}{\small\color{kfInk} 위에서 배운 낱말을 넣어 빈칸을 채워 보세요.}\par\nopagebreak[4]\smallskip\nopagebreak[4]
\begin{kfBlankList}
\kfBlankItem 나의 \kfFillBlank[12mm]{}은/는 밝고 활발합니다.
\kfBlankItem 나의 \kfFillBlank[12mm]{}은/는 그림 그리기입니다.
\kfBlankItem 나의 \kfFillBlank[12mm]{}은/는 친구들을 잘 웃기는 것입니다.
\kfBlankItem 나의 \kfFillBlank[12mm]{}은/는 훌륭한 과학자가 되는 것입니다.
\end{kfBlankList}

\kfActivityBg \par\kfMaybeNeedspace{24\baselineskip}\noindent{\kfTipMark\,\,}{\small\color{kfInk} 문제를 읽고 O, X 표시를 해 보세요.}\par\nopagebreak[4]\smallskip\nopagebreak[4]
\begin{kfOXList}
\kfOXItem 자기소개를 할 때는 거짓말을 해도 괜찮다. \kfFillBlank[12mm]{} \kfOX
\kfOXItem 구체적으로 말하면 사람들이 나를 더 잘 이해할 수 있다. \kfFillBlank[12mm]{} \kfOX
\kfOXItem 자기소개를 할 때는 주로 부정적인 말을 해야 한다. \kfFillBlank[12mm]{} \kfOX
\kfOXItem "나는 매주 토요일마다 축구를 해요"는 구체적인 표현이다. \kfFillBlank[12mm]{} \kfOX
\kfOXItem 좋은 자기소개는 다른 사람에게 나를 잘 알릴 수 있다. \kfFillBlank[12mm]{} \kfOX
\end{kfOXList}

\kfSecHeader{kfAreaConcept}{문제}{}

\begin{kfQBlock}{01}{객관식}
\kfQStem{다음 중 올바른 문장 구조를 가진 것은?}
\begin{kfChoices}
\kfChoice 예쁜 꽃이
\kfChoice 아이가 웃습니다
\kfChoice 빨리 달립니다
\kfChoice 맛있는
\end{kfChoices}
\end{kfQBlock}
\begin{kfQBlock}{02}{서술형}
\kfQStem{다음 문장에서 주어와 서술어를 찾아보세요.}
\kfAnswerNote{답안}{4}
\end{kfQBlock}


\kfStartLit{「나는 나야, 특별한 나」}


\kfPassageNote{}{%
나는 나야, 특별한 나 \\[4pt]
세상에 하나뿐인 소중한 나 \\[4pt]
까만 머리, 반짝이는 눈 \\[4pt]
환하게 웃는 내 얼굴 \\[4pt]
거울 속에 비친 나를 봐 \\[4pt]
책 읽기를 좋아하는 나 \\[4pt]
친구들과 뛰노는 걸 좋아하는 나 \\[4pt]
엄마 아빠 사랑하는 착한 나 \\[4pt]
때로는 실수도 하지만 \\[4pt]
다시 일어서는 씩씩한 나 \\[4pt]
나는 나야, 멋진 나!
}


\kfActivityBg \par\kfMaybeNeedspace{18\baselineskip}\noindent{\kfTipMark\,\,}{\small\color{kfInk} 주어진 안내에 맞춰 자유롭게 글로 표현해 보세요.}\par\nopagebreak[4]\smallskip\nopagebreak[4]
\kfWritingPrompt{'거울 속의 나'를 읽고, 나의 좋은 점을 생각하며 나를 자랑하는 짧은 글을 자유롭게 써보세요.}
\kfWriting{답안 작성}{8}

\kfSecHeader{kfAreaLiterature}{문제}{}

\begin{kfQBlock}{01}{객관식}
\kfQStem{시에서 '나'가 좋아하는 것 두 가지는 무엇인가요?}
\begin{kfChoices}
\kfChoice 게임하기와 TV 보기
\kfChoice 책 읽기와 친구들과 놀기
\kfChoice 그림 그리기와 노래 부르기
\kfChoice 운동하기와 요리하기
\end{kfChoices}
\end{kfQBlock}
\begin{kfQBlock}{02}{객관식}
\kfQStem{시에서 '나'는 자신을 어떻게 표현했나요?}
\begin{kfChoices}
\kfChoice 평범하고 특별할 것 없는 사람
\kfChoice 특별하고 소중한 사람
\kfChoice 완벽하고 실수하지 않는 사람
\kfChoice 다른 사람보다 못한 사람
\end{kfChoices}
\end{kfQBlock}
\begin{kfQBlock}{03}{객관식}
\kfQStem{이 시에서 '나'의 모습으로 알맞지 \uline{\underline{않은}} 것은 무엇인가요?}
\begin{kfChoices}
\kfChoice 까만 머리에 반짝이는 눈을 가졌다.
\kfChoice 거울 속에 비친 모습이 환하게 웃는다.
\kfChoice 실수를 한 번도 하지 않는 완벽한 모습이다.
\kfChoice 실수해도 다시 일어서는 씩씩한 모습이다.
\end{kfChoices}
\end{kfQBlock}
\begin{kfQBlock}{04}{서술형}
\kfQStem{이 시에서 '나'의 외모를 표현한 부분을 모두 찾아 쓰세요.}
\kfAnswerNote{답안}{4}
\end{kfQBlock}
\begin{kfQBlock}{05}{객관식}
\kfQStem{"세상에 하나뿐인 소중한 나"라는 말에서 알 수 있는 것은?}
\begin{kfChoices}
\kfChoice 모든 사람은 똑같다.
\kfChoice 나는 특별하지 않다.
\kfChoice 모든 사람은 각자 특별하다.
\kfChoice 나보다 중요한 사람이 많다.
\end{kfChoices}
\end{kfQBlock}
\begin{kfQBlock}{06}{객관식}
\kfQStem{이 시의 마지막 행 "나는 나야, 멋진 나!"에 가장 어울리는 시의 주제는 무엇인가요?}
\begin{kfChoices}
\kfChoice 세상에 하나뿐인 나를 소중히 여기자.
\kfChoice 책 읽기와 뛰노는 일 가운데 한 가지만 잘하자.
\kfChoice 실수가 적은 친구를 따라 똑같이 행동하자.
\kfChoice 까만 머리와 반짝이는 눈을 가꾸어 멋져 보이자.
\end{kfChoices}
\end{kfQBlock}
\begin{kfQBlock}{07}{서술형}
\kfQStem{이 시에서 '나'가 자신을 긍정적으로 생각하는 이유는 무엇일까요?}
\kfAnswerNote{답안}{4}
\end{kfQBlock}


\kfStartNonlit{나를 알리는 자기소개 쓰기}


\kfPassageNote{나를 알리는 자기소개 쓰기}{%
나를 처음 만나는 사람들에게 나에 대해 알려주는 글을 자기소개 글이라고 합니다. 자기소개 글을 쓰기 위해서는 가장 먼저 첫 문장에 자신의 이름을 정확하게 써야 합니다. "안녕하세요. 저는 김민수입니다."와 같이 정중한 인사와 함께 이름을 밝히는 것이 예의 있는 시작입니다.  이름을 말한 뒤에는 나의 구체적인 특징을 소개합니다. 내가 몇 살인지, 몇 학년인지, 그리고 어디에 살고 있는지 등을 차례대로 적습니다. 그다음에는 내가 평소에 무엇을 좋아하는지 써 봅니다. 단순히 "운동을 좋아합니다."라고 하기보다는 "저는 축구를 좋아합니다."처럼 구체적인 활동을 적는 것이 좋습니다.  글의 마지막 부분에는 나의 장래 희망이나 목표를 담습니다. 나중에 커서 어떤 사람이 되고 싶은지 꿈을 밝히면 나를 더 인상 깊게 알릴 수 있습니다. 모든 내용을 다 적었다면 "앞으로 잘 부탁드립니다."라는 마무리 인사로 글을 끝맺습니다. 정해진 순서에 따라 차근차근 글을 쓰면 누구나 멋진 자기소개 글을 완성할 수 있습니다.
}


\kfActivityBg \par\needspace{15\baselineskip}
\kfSecHeader{kfAreaNonlit}{활동}{문장 독해}
\par\kfMaybeNeedspace{32\baselineskip}\noindent{\kfTipMark\,\,}{\small\color{kfInk} 각 문장을 읽고 물음에 답하세요.}\par\nopagebreak[4]\smallskip\nopagebreak[4]
\begin{kfSentenceUnit}{1}{좋아합니다. 축구를 저는}
\kfSentQ{1.}{(1) 이 문장에서 '누가'에 해당하는 것은 무엇인가요?}
\begin{kfChoices}
\kfChoice 좋아합니다.
\kfChoice 축구를
\kfChoice 저는
\end{kfChoices}
\kfSentQ{1.}{(2) 이 문장에서 '어떠하다'에 해당하는 것은 무엇인가요?}
\begin{kfChoices}
\kfChoice 축구를
\kfChoice 좋아합니다
\kfChoice 저는
\end{kfChoices}
\end{kfSentenceUnit}

\begin{kfSentenceUnit}{2}{멋진 자기소개 있습니다. 누구나}
\kfSentQ{2.}{(1) 이 문장에서 '누가'에 해당하는 것은 무엇인가요?}
\begin{kfChoices}
\kfChoice 멋진 자기소개
\kfChoice 있습니다.
\kfChoice 누구나
\end{kfChoices}
\kfSentQ{2.}{(2) 이 문장에서 '어찌하다'에 해당하는 것은 무엇인가요?}
\begin{kfChoices}
\kfChoice 자기소개 글을
\kfChoice 완성할 수 있습니다
\kfChoice 누구나
\end{kfChoices}
\end{kfSentenceUnit}

\kfActivityBg \par\kfMaybeNeedspace{24\baselineskip}\noindent{\kfTipMark\,\,}{\small\color{kfInk} 빈칸에 알맞은 말을 채워 보세요.}\par\nopagebreak[4]\smallskip\nopagebreak[4]
\begin{kfBlankList}
\kfBlankItem 글을 읽고 빈칸에 알맞은 말을 넣어 내용을 요약해 봅시다.

자기소개 글을 쓸 때는 정해진 순서가 있습니다.  먼저 첫 문장에 인사와 함께 자신의 ①\kfFillBlank[12mm]{}을 밝힙니다.   그다음 나이나 사는 곳 같은 특징을 소개하고, 내가 좋아하는 것을 ②\kfFillBlank[12mm]{}으로 적습니다.  마지막에는 나의 ③\kfFillBlank[12mm]{}이나 목표를 쓰고 정중한 인사를 하며 글을 마칩니다.
\end{kfBlankList}

\kfActivityBg \par\kfMaybeNeedspace{18\baselineskip}\noindent{\kfTipMark\,\,}{\small\color{kfInk} 주어진 안내에 맞춰 자유롭게 글로 표현해 보세요.}\par\nopagebreak[4]\smallskip\nopagebreak[4]
\kfWritingPrompt{배운 순서에 따라 나를 소개하는 글을 써 봅시다.}
\kfWriting{답안 작성}{8}

\kfSecHeader{kfAreaNonlit}{문제}{}

\begin{kfQBlock}{01}{객관식}
\kfQStem{이 글에서 자기소개 글의 첫 문장의 예로 든 것은 무엇인가요?}
\begin{kfChoices}
\kfChoice "저는 축구를 좋아합니다."
\kfChoice "안녕하세요. 저는 김민수입니다."
\kfChoice "커서 멋진 사람이 되고 싶습니다."
\kfChoice "앞으로 잘 부탁드립니다."
\end{kfChoices}
\end{kfQBlock}
\begin{kfQBlock}{02}{객관식}
\kfQStem{이 글에서 이름을 말한 뒤 소개하라고 한 '나의 구체적인 특징'에 들어가지 \uline{않는} 것은 무엇인가요?}
\begin{kfChoices}
\kfChoice 올해 내가 정확히 몇 살이 되었는지
\kfChoice 내가 다니는 학교의 몇 학년인지
\kfChoice 내가 지금 어떤 동네에 살고 있는지
\kfChoice 내가 글의 마지막에 쓸 마무리 인사
\end{kfChoices}
\end{kfQBlock}
\begin{kfQBlock}{03}{객관식}
\kfQStem{좋아하는 것을 쓸 때 더 좋은 방법은 무엇인가요?}
\begin{kfChoices}
\kfChoice 좋아하는 이유만 길게 씁니다.
\kfChoice 좋아하는 것을 하나도 쓰지 않습니다.
\kfChoice 축구처럼 구체적인 활동을 적습니다.
\kfChoice 단순히 운동이라고만 짧게 말합니다.
\end{kfChoices}
\end{kfQBlock}
\begin{kfQBlock}{04}{객관식}
\kfQStem{이 글에서 자기소개 글의 마지막 부분에 담으라고 한 내용은 무엇인가요?}
\begin{kfChoices}
\kfChoice 나의 장래 희망이나 목표
\kfChoice 좋아하는 활동을 구체적으로 적은 한 문장
\kfChoice 자신의 이름과 나이
\kfChoice 사는 곳과 학년
\end{kfChoices}
\end{kfQBlock}
\begin{kfQBlock}{05}{객관식}
\kfQStem{이 글에서 자기소개 글을 끝맺는 마무리 인사로 알려 준 표현은 무엇인가요?}
\begin{kfChoices}
\kfChoice "안녕하세요. 저는 김민수입니다."
\kfChoice "저는 축구를 좋아합니다."
\kfChoice "나중에 커서 의사가 되고 싶습니다."
\kfChoice "앞으로 잘 부탁드립니다."
\end{kfChoices}
\end{kfQBlock}


\kfStartWeekly{실력확인 영역 학습}


\noindent\textit{\small 제한시간: 20분 / 총 10문항}\par\medskip
\par\noindent\textbf{[4\textasciitilde{}5] 다음 글을 읽고 물음에 답하시오.}\par\medskip
\kfPassageNote{지문 4}{%
나는 나야, 특별한 나 \\[4pt]
세상에 하나뿐인 소중한 나 \\[4pt]
까만 머리, 반짝이는 눈 \\[4pt]
환하게 웃는 내 얼굴 \\[4pt]
거울 속에 비친 나를 봐 \\[4pt]
책 읽기를 좋아하는 나 \\[4pt]
친구들과 뛰노는 걸 좋아하는 나 \\[4pt]
엄마 아빠 사랑하는 착한 나 \\[4pt]
때로는 실수도 하지만 \\[4pt]
다시 일어서는 씩씩한 나 \\[4pt]
나는 나야, 멋진 나!
}
\par\noindent\textbf{[6\textasciitilde{}7] 다음 글을 읽고 물음에 답하시오.}\par\medskip
\kfPassageNote{지문 6}{%
나를 처음 만나는 사람들에게 나에 대해 알려주는 글을 자기소개 글이라고 합니다. 자기소개 글을 쓰기 위해서는 가장 먼저 첫 문장에 자신의 이름을 정확하게 써야 합니다. "안녕하세요. 저는 김민수입니다."와 같이 정중한 인사와 함께 이름을 밝히는 것이 예의 있는 시작입니다.  이름을 말한 뒤에는 나의 구체적인 특징을 소개합니다. 내가 몇 살인지, 몇 학년인지, 그리고 어디에 살고 있는지 등을 차례대로 적습니다. 그다음에는 내가 평소에 무엇을 좋아하는지 써 봅니다. 단순히 "운동을 좋아합니다."라고 하기보다는 "저는 축구를 좋아합니다.
}
\begin{kfQBlock}{01}{객관식}
\kfQStem{'특기'의 뜻은 무엇인가요?}
\begin{kfChoices}
\kfChoice 특별한 날
\kfChoice 특별한 기분
\kfChoice 특별한 장소
\kfChoice 남보다 특별히 잘하는 것
\end{kfChoices}
\end{kfQBlock}

\begin{kfQBlock}{02}{객관식}
\kfQStem{다음 중 낱말의 쓰임이 \uline{어색한} 문장은 무엇인가요?}
\begin{kfChoices}
\kfChoice 제 취미는 축구입니다.
\kfChoice 제 꿈은 요리사입니다.
\kfChoice 온 가족이 모였습니다.
\kfChoice 제 성격은 피아노입니다.
\end{kfChoices}
\end{kfQBlock}

\begin{kfQBlock}{03}{객관식}
\kfQStem{자기소개를 할 때 알맞지 \uline{\underline{않은}} 태도는?}
\begin{kfChoices}
\kfChoice 자신의 이름을 가장 먼저 밝힌다
\kfChoice 좋아하는 것과 잘하는 것을 함께 말한다
\kfChoice 나를 돋보이게 하려고 내용을 부풀려서 꾸민다
\kfChoice 듣는 사람의 눈을 보며 또박또박 말한다
\end{kfChoices}
\end{kfQBlock}

\begin{kfQBlock}{04}{객관식}
\kfQStem{시에서 '나'가 실수했을 때 보여주는 모습은 무엇인가요?}
\begin{kfChoices}
\kfChoice 거울 속 자기 모습을 한참 들여다본다
\kfChoice 씩씩하게 일어선다
\kfChoice 엄마 아빠께 도와 달라고 부탁한다
\kfChoice 책을 읽으며 마음을 가라앉힌다
\end{kfChoices}
\end{kfQBlock}

\begin{kfQBlock}{05}{객관식}
\kfQStem{시에서 알 수 있는 글쓴이의 생각은 무엇인가요?}
\begin{kfChoices}
\kfChoice 나는 특별하고 소중하다
\kfChoice 나는 공부를 아주 잘한다
\kfChoice 나는 실수를 전혀 안 한다
\kfChoice 나는 세상에서 제일 예쁘다
\end{kfChoices}
\end{kfQBlock}

\begin{kfQBlock}{06}{객관식}
\kfQStem{'자기소개 글'의 특징으로 알맞은 것은 무엇인가요?}
\begin{kfChoices}
\kfChoice 인사는 생략해도 됩니다.
\kfChoice 정해진 순서에 따라 씁니다.
\kfChoice 글의 첫머리에 꿈을 적습니다.
\kfChoice 나를 돋보이게 하려고 내용을 거짓으로 씁니다.
\end{kfChoices}
\end{kfQBlock}

\begin{kfQBlock}{07}{객관식}
\kfQStem{자기소개 글의 첫 문장에 반드시 포함되어야 할 내용은 무엇인가요?}
\begin{kfChoices}
\kfChoice 나의 이름
\kfChoice 내일의 날씨
\kfChoice 어제 본 영화
\kfChoice 좋아하는 음식
\end{kfChoices}
\end{kfQBlock}

\begin{kfQBlock}{08}{객관식}
\kfQStem{문장을 이룰 때 '누가' 혹은 '무엇이'에 해당하는 말은 무엇인가요?}
\begin{kfChoices}
\kfChoice 제목
\kfChoice 특징
\kfChoice 주어
\kfChoice 서술어
\end{kfChoices}
\end{kfQBlock}

\begin{kfQBlock}{09}{서술형}
\kfQStem{다음 문장에서 '누가'에 해당하는 것을 찾아 쓰세요.}
\begin{kfEvidence}"민수는 노래를 부릅니다."\end{kfEvidence}
\kfAnswerNote{답안}{4}
\end{kfQBlock}

\begin{kfQBlock}{10}{서술형}
\kfQStem{다음 문장을 '누가'와 '무엇이다'에 해당하는 것을 나누어 쓰세요.}
\begin{kfEvidence}"나의 꿈은 경찰관입니다."\end{kfEvidence}
\kfAnswerNote{답안}{4}
\end{kfQBlock}




\clearpage

\kfChapterCover{2}{소쉬르2 (20주차) (2)}{Chapter 2}{이번 챕터의 학습 내용을 시작합니다.}
\begin{kfChapterAreaList}
\kfChapterAreaItem{어휘}{kfAreaVocab}{마음을 표현하는 낱말 익히기}
\kfChapterAreaItem{개념}{kfAreaConcept}{감정은 날씨와 같아요}
\kfChapterAreaItem{문학}{kfAreaLiterature}{「내 마음은 알록달록 도화지.」}
\kfChapterAreaItem{비문학}{kfAreaNonlit}{화를 다스리는 마음의 힘}
\kfChapterAreaItem{실력확인}{kfAreaWeekly}{}
\end{kfChapterAreaList}

\kfStartVocab{마음을 표현하는 낱말 익히기}


\begin{kfVocab}
\kfVocabItem{1}{욕구가 채워지거나 좋은 일이 생겨서 즐거운 마음}
\kfVocabItem{2}{겁이 나고 무서운 마음}
\kfVocabItem{3}{스르르 고개를 숙이게 되는 수줍은 마음}
\kfVocabItem{4}{벅찰 만큼 기특하고 자랑스러운 마음}
\kfVocabItem{5}{마음이 아프고 눈물이 날 것 같은 느낌}
\kfVocabItem{6}{기대했던 일이 이루어지지 않아서 기운이 빠지는 마음}
\end{kfVocab}


\kfStartConcept{감정은 날씨와 같아요}


\kfPassageNoteNT{%
우리 마음속 감정은 마치 하늘의 날씨와 비슷해요.

맑은 날이 있으면 흐린 날도 있듯이, 감정이 변하는 것은 자연스러운 일이에요.

그리고 '화'나 '슬픔'도 우리에게 꼭 필요한 감정이니 슬플 때는 실컷 울어도 괜찮아요.

하지만 비가 많이 오면 우산을 쓰듯이 화가 많이 날 때는 마음을 잘 다스려야 한답니다.
\par\medskip\begin{itemize}[leftmargin=18pt, itemsep=2pt, topsep=2pt, label=\textbullet]
\item 맑은 날 — 기쁨
\item 흐린 날 — 슬픔
\item 비 오는 날 — 화 (마음 다스리기 필요)
\end{itemize}
}

\kfPassageNote{문장 독해와 감정 표현}{%
문장에서 사물이나 사람의 상태, 그리고 마음의 느낌을 설명하는 말을 라고 해요. 우리말에는 감정을 나타내는 형용사가 아주 많아요. 예시: 기쁘다, 슬프다, 무섭다, 즐겁다, 외롭다
}


\kfActivityBg \par\kfMaybeNeedspace{24\baselineskip}\noindent{\kfTipMark\,\,}{\small\color{kfInk} 위에서 배운 낱말 중 알맞은 것을 골라 빈칸을 채워 보세요.}\par\nopagebreak[4]\smallskip\nopagebreak[4]
\begin{kfBlankList}
\kfBlankItem 칭찬을 받으니 마음이 \kfFillBlank[12mm]{}합니다.
\kfBlankItem 아끼던 인형을 잃어버려서 \kfFillBlank[12mm]{}에 잠겼습니다.
\kfBlankItem 사람들 앞에 서니 얼굴이 빨개지고 \kfFillBlank[12mm]{}이 듭니다.
\kfBlankItem 친구가 장난감을 망가뜨려 놓고도 미안해하지 않아서 \kfFillBlank[12mm]{}했습니다.
\end{kfBlankList}

\kfActivityBg \par\kfMaybeNeedspace{24\baselineskip}\noindent{\kfTipMark\,\,}{\small\color{kfInk} 문제를 읽고 O, X 표시를 해 보세요.}\par\nopagebreak[4]\smallskip\nopagebreak[4]
\begin{kfOXList}
\kfOXItem 우리는 다양한 감정을 느낄 수 있다. \kfFillBlank[12mm]{} \kfOX
\kfOXItem 감정은 날씨처럼 계속 변할 수 있다. \kfFillBlank[12mm]{} \kfOX
\kfOXItem 화가 날 때는 무조건 참기만 해야 한다. \kfFillBlank[12mm]{} \kfOX
\kfOXItem 슬프거나 화가 나는 것은 무조건 나쁜 것이다. \kfFillBlank[12mm]{} \kfOX
\kfOXItem 감정을 잘 다스리는 것은 우산을 쓰는 것과 비슷하다. \kfFillBlank[12mm]{} \kfOX
\end{kfOXList}

\kfActivityBg \par\kfMaybeNeedspace{24\baselineskip}\noindent{\kfTipMark\,\,}{\small\color{kfInk} 다음 활동을 풀어 보세요.}\par\nopagebreak[4]\smallskip\nopagebreak[4]
\begin{kfConnect}
\kfPair{다음 상황에 어울리는 감정 낱말을 찾아 연결해 보세요.}{}
\end{kfConnect}

\kfSecHeader{kfAreaConcept}{문제}{}

\begin{kfQBlock}{01}{객관식}
\kfQStem{다음 중 감정을 나타내는 말이 \underline{아닌} 것은?}
\begin{kfChoices}
\kfChoice 즐겁다
\kfChoice 달린다
\kfChoice 슬프다
\kfChoice 무섭다
\end{kfChoices}
\end{kfQBlock}


\kfStartLit{「내 마음은 알록달록 도화지.」}


\kfPassageNote{}{%
내 마음은 알록달록 도화지. \\[4pt]
맛있는 아이스크림 먹을 땐 \\[4pt]
노랑노랑 기쁨의 색. \\[4pt]
넘어져서 무릎이 까진 날엔 \\[4pt]
파랑파랑 슬픔의 색. \\[4pt]
동생이 내 그림에 낙서하면 \\[4pt]
빨강빨강 화나는 색. \\[4pt]
오늘은 무슨 색을 칠해 볼까? \\[4pt]
어떤 색이라도 괜찮아. \\[4pt]
모두 다 소중한 내 마음이니까.
}


\kfActivityBg \par\kfMaybeNeedspace{18\baselineskip}\noindent{\kfTipMark\,\,}{\small\color{kfInk} 주어진 안내에 맞춰 자유롭게 글로 표현해 보세요.}\par\nopagebreak[4]\smallskip\nopagebreak[4]
\kfWritingPrompt{오늘 내 마음의 날씨는 어떤가요? 맑음, 흐림, 비, 천둥번개 중에서 골라 이유와 함께 짧게 써 보세요.}
\kfWriting{답안 작성}{8}

\kfSecHeader{kfAreaLiterature}{문제}{}

\begin{kfQBlock}{01}{객관식}
\kfQStem{시에서 '나'가 슬픔을 느낀 이유는 무엇인가요?}
\begin{kfChoices}
\kfChoice 차가운 아이스크림을 먹어서
\kfChoice 넘어져서 무릎이 까져서
\kfChoice 동생이 책에 낙서를 해서
\kfChoice 풀어야 할 숙제가 너무 많아서
\end{kfChoices}
\end{kfQBlock}
\begin{kfQBlock}{02}{객관식}
\kfQStem{시에서 '나'가 동생 때문에 느낀 감정은 무엇인가요?}
\begin{kfChoices}
\kfChoice 기쁨
\kfChoice 슬픔
\kfChoice 화
\kfChoice 뿌듯함
\end{kfChoices}
\end{kfQBlock}
\begin{kfQBlock}{03}{객관식}
\kfQStem{"어떤 색이라도 괜찮아"라는 말의 뜻으로 알맞은 것은?}
\begin{kfChoices}
\kfChoice 빨간색만 칠해야 한다는 뜻이다.
\kfChoice 어떤 감정이든 모두 소중하다는 뜻이다.
\kfChoice 그림을 그릴 때는 아무 색이나 써야 한다는 뜻이다.
\kfChoice 기쁜 마음만 가져야 한다는 뜻이다.
\end{kfChoices}
\end{kfQBlock}
\begin{kfQBlock}{04}{서술형}
\kfQStem{이 시에서 '기쁨'을 무슨 색으로 표현했나요?}
\kfAnswerNote{답안}{4}
\end{kfQBlock}
\begin{kfQBlock}{05}{객관식}
\kfQStem{이 시의 제목을 다른 것으로 바꾼다면 무엇이 좋을까요?}
\begin{kfChoices}
\kfChoice 맛있는 아이스크림
\kfChoice 미운 내 동생
\kfChoice 내 마음속 여러 가지 색깔
\kfChoice 그림 그리기 대회
\end{kfChoices}
\end{kfQBlock}
\begin{kfQBlock}{06}{객관식}
\kfQStem{이 시를 읽고 알 수 있는 내용이 \uline{\underline{아닌}} 것은?}
\begin{kfChoices}
\kfChoice 감정은 여러 가지가 있다.
\kfChoice 상황에 따라 마음이 달라진다.
\kfChoice 화가 나는 마음은 숨겨야 한다.
\kfChoice 내 마음을 솔직하게 표현해도 된다.
\end{kfChoices}
\end{kfQBlock}
\begin{kfQBlock}{07}{서술형}
\kfQStem{시의 마지막 부분에서 "모두 다 소중한 내 마음이니까"라고 말한 이유는 무엇일까요?}
\kfAnswerNote{답안}{4}
\end{kfQBlock}


\kfStartNonlit{화를 다스리는 마음의 힘}


\kfPassageNote{화를 다스리는 마음의 힘}{%
우리는 생활하면서 화가 나거나 속상한 마음이 들 때가 있습니다. 누군가 내 장난감을 망가뜨리거나 나를 놀리면 기분이 나빠지기 마련입니다. 하지만 화가 난다고 해서 소리를 지르거나 화풀이를 하면 기분이 더 안 좋아질 수 있습니다. 이럴 때는 마음을 지키는 방법을 사용해 보아야 합니다. 가장 먼저 마음속으로 숫자를 천천히 세어 보세요. 화가 머리 끝까지 날 때 1부터 10까지 숫자를 세면 화가 난 마음이 조금씩 차분해집니다. 그다음에는 풍선이 부풀어 오르는 것처럼 배를 움직이며 숨을 쉬어 보세요. 코로 숨을 깊게 들이마시고 입으로 천천히 내쉬면 몸의 긴장이 풀리면서 마음이 편안해집니다. 마지막으로 내 마음을 솔직하게 말하는 연습이 필요합니다. 친구에게 화를 내는 대신 내 기분을 솔직하게 전해 보세요. "너 때문에 망가졌잖아!"라고 소리를 지르기보다 "네가 장난감을 망가뜨려서 나는 정말 속상해."라고 내 기분을 전하는 것입니다. 이렇게 말하면 상대방도 내 마음을 더 잘 이해할 수 있습니다.
}


\kfActivityBg \par\needspace{15\baselineskip}
\kfSecHeader{kfAreaNonlit}{활동}{문장 독해}
\par\kfMaybeNeedspace{32\baselineskip}\noindent{\kfTipMark\,\,}{\small\color{kfInk} 각 문장을 읽고 물음에 답하세요.}\par\nopagebreak[4]\smallskip\nopagebreak[4]
\begin{kfSentenceUnit}{1}{민수는 숫자를 마음속으로}
\kfSentQ{1.}{(1) 이 문장에서 '누가'에 해당하는 것은 무엇인가요?}
\begin{kfChoices}
\kfChoice 민수는
\kfChoice 숫자를
\kfChoice 마음속으로
\end{kfChoices}
\kfSentQ{1.}{(2) 이 문장에서 '무엇을'에 해당하는 것은 무엇인가요?}
\begin{kfChoices}
\kfChoice 마음속으로
\kfChoice 센다
\kfChoice 숫자를
\end{kfChoices}
\kfSentQ{1.}{(3) 이 문장에서 '어찌하다'에 해당하는 것은 무엇인가요?}
\begin{kfChoices}
\kfChoice 센다
\kfChoice 마음속으로
\kfChoice 숫자를
\end{kfChoices}
\end{kfSentenceUnit}

\begin{kfSentenceUnit}{2}{상대방은 이해한다 마음을}
\kfSentQ{2.}{(1) 이 문장에서 '누가'에 해당하는 것은 무엇인가요?}
\begin{kfChoices}
\kfChoice 상대방은
\kfChoice 이해한다
\kfChoice 마음을
\end{kfChoices}
\kfSentQ{2.}{(2) 이 문장에서 '무엇을'에 해당하는 것은 무엇인가요?}
\begin{kfChoices}
\kfChoice 상대방은
\kfChoice 내 마음을
\kfChoice 이해한다
\end{kfChoices}
\kfSentQ{2.}{(3) 이 문장에서 '어찌하다'에 해당하는 것은 무엇인가요?}
\begin{kfChoices}
\kfChoice 이해한다
\kfChoice 상대방은
\kfChoice 내 마음을
\end{kfChoices}
\end{kfSentenceUnit}

\kfActivityBg \par\kfMaybeNeedspace{24\baselineskip}\noindent{\kfTipMark\,\,}{\small\color{kfInk} 빈칸에 알맞은 말을 채워 보세요.}\par\nopagebreak[4]\smallskip\nopagebreak[4]
\begin{kfBlankList}
\kfBlankItem 글을 읽고 빈칸에 알맞은 말을 넣어 내용을 정리해 봅시다.

화가 날 때는 마음을 지키는 방법을 사용해야 합니다.  먼저 마음속으로 ①\kfFillBlank[12mm]{}를 천천히 세면 마음이 차분해집니다.  그다음에는 풍선처럼 배를 움직이며 깊게 ②\kfFillBlank[12mm]{}를 쉽니다.  마지막으로 친구에게 소리를 지르는 대신 나의 ③\kfFillBlank[12mm]{}을 솔직하게 말하는 연습을 합니다.
\end{kfBlankList}

\kfActivityBg \par\kfMaybeNeedspace{18\baselineskip}\noindent{\kfTipMark\,\,}{\small\color{kfInk} 주어진 안내에 맞춰 자유롭게 글로 표현해 보세요.}\par\nopagebreak[4]\smallskip\nopagebreak[4]
\kfWritingPrompt{다음 상황에서 내 기분을 솔직하게 말하는 문장을 완성해 보세요.}
\kfWriting{답안 작성}{8}

\kfSecHeader{kfAreaNonlit}{문제}{}

\begin{kfQBlock}{01}{객관식}
\kfQStem{이 글에서 화가 난다고 해서 소리를 지르거나 화풀이를 하면 어떻게 된다고 했나요?}
\begin{kfChoices}
\kfChoice 기분이 더 안 좋아질 수 있다.
\kfChoice 화가 머리 끝까지 났을 때 1부터 10까지 빨리 셀 수 있다.
\kfChoice 풍선처럼 배를 움직이며 깊은 숨이 저절로 쉬어진다.
\kfChoice 코로 들이마신 숨이 입으로 천천히 빠져나간다.
\end{kfChoices}
\end{kfQBlock}
\begin{kfQBlock}{02}{객관식}
\kfQStem{이 글에서 화가 났을 때 '가장 먼저' 해 보라고 한 방법은 무엇인가요?}
\begin{kfChoices}
\kfChoice 친구에게 내 기분을 솔직하게 말한다.
\kfChoice 마음속으로 1부터 10까지 숫자를 천천히 센다.
\kfChoice 풍선처럼 배를 움직이며 숨을 쉰다.
\kfChoice 장난감을 망가뜨린 친구에게 큰 소리로 외친다.
\end{kfChoices}
\end{kfQBlock}
\begin{kfQBlock}{03}{객관식}
\kfQStem{이 글에서 코로 깊게 들이마시고 입으로 천천히 내쉴 때 몸에 일어난다고 한 변화는 무엇인가요?}
\begin{kfChoices}
\kfChoice 화가 머리 끝까지 솟구친다.
\kfChoice 몸의 긴장이 풀리면서 마음이 편안해진다.
\kfChoice 마음속 숫자가 1부터 10까지 저절로 세어진다.
\kfChoice 망가진 장난감이 다시 처음 모양으로 돌아온다.
\end{kfChoices}
\end{kfQBlock}
\begin{kfQBlock}{04}{객관식}
\kfQStem{이 글에서 친구가 장난감을 망가뜨렸을 때, 화를 내는 대신 권한 표현은 무엇인가요?}
\begin{kfChoices}
\kfChoice "너 때문에 내 소중한 장난감이 다 망가졌잖아!"라고 큰 소리로 외친다.
\kfChoice 마음속으로 1부터 10까지 숫자를 천천히 세어 화난 마음을 차분히 가라앉힌다.
\kfChoice 코로 숨을 깊게 들이마시고 입으로 천천히 내쉬며 몸의 긴장을 풀어 준다.
\kfChoice "네가 장난감을 망가뜨려서 나는 속상해."라고 솔직히 전한다.
\end{kfChoices}
\end{kfQBlock}
\begin{kfQBlock}{05}{객관식}
\kfQStem{이 글에서 내 기분을 솔직하게 말했을 때 생긴다고 한 좋은 점은 무엇인가요?}
\begin{kfChoices}
\kfChoice 망가진 장난감이 새것이 된다.
\kfChoice 화가 났던 일이 즉시 모두 사라진다.
\kfChoice 상대방도 내 마음을 더 잘 이해할 수 있다.
\kfChoice 마음속 숫자를 더 빨리 셀 수 있게 된다.
\end{kfChoices}
\end{kfQBlock}


\kfStartWeekly{실력확인 영역 학습}


\noindent\textit{\small 제한시간: 20분 / 총 10문항}\par\medskip
\par\noindent\textbf{[4\textasciitilde{}5] 다음 글을 읽고 물음에 답하시오.}\par\medskip
\kfPassageNote{지문 4}{%
내 마음은 알록달록 도화지. \\[4pt]
맛있는 아이스크림 먹을 땐 \\[4pt]
노랑노랑 기쁨의 색. \\[4pt]
넘어져서 무릎이 까진 날엔 \\[4pt]
파랑파랑 슬픔의 색. \\[4pt]
동생이 내 그림에 낙서하면 \\[4pt]
빨강빨강 화나는 색. \\[4pt]
오늘은 무슨 색을 칠해 볼까? \\[4pt]
어떤 색이라도 괜찮아. \\[4pt]
모두 다 소중한 내 마음이니까.
}
\par\noindent\textbf{[6\textasciitilde{}8] 다음 글을 읽고 물음에 답하시오.}\par\medskip
\kfPassageNote{지문 6}{%
우리는 생활하면서 화가 나거나 속상한 마음이 들 때가 있습니다. 누군가 내 장난감을 망가뜨리거나 나를 놀리면 기분이 나빠지기 마련입니다. 하지만 화가 난다고 해서 소리를 지르거나 화풀이를 하면 기분이 더 안 좋아질 수 있습니다. 이럴 때는 마음을 지키는 방법을 사용해 보아야 합니다. 가장 먼저 마음속으로 숫자를 천천히 세어 보세요. 화가 머리 끝까지 날 때 1부터 10까지 숫자를 세면 화가 난 마음이 조금씩 차분해집니다. 그다음에는 풍선이 부풀어 오르는 것처럼 배를 움직이며 숨을 쉬어 보세요. 코로 숨을 깊게 들이마시고 입으로 천천히 내쉬면 몸의 긴장이 풀리면서 마음이 편안해집니다. 마지막으로 내 마음을 솔직하게 말하는 연습이 필요합니다. 친구에게 화를 내는 대신 내 기분을 솔직하게 전해 보세요. "너 때문에 망가졌잖아!"라고 소리를 지르기보다 "네가 장난감을 망가뜨려서 나는 정말 속상해."라고 내 기분을 전하는 것입니다. 이렇게 말하면 상대방도 내 마음을 더 잘 이해할 수 있습니다.
}
\begin{kfQBlock}{01}{객관식}
\kfQStem{욕구가 채워지거나 좋은 일이 생겨서 즐거운 마음을 무엇이라고 하나요?}
\begin{kfChoices}
\kfChoice 화
\kfChoice 슬픔
\kfChoice 기쁨
\kfChoice 두려움
\end{kfChoices}
\end{kfQBlock}

\begin{kfQBlock}{02}{객관식}
\kfQStem{'두려움'과 비슷한 뜻을 가진 낱말은?}
\begin{kfChoices}
\kfChoice 고마움
\kfChoice 즐거움
\kfChoice 무서움
\kfChoice 반가움
\end{kfChoices}
\end{kfQBlock}

\begin{kfQBlock}{03}{객관식}
\kfQStem{마음속 감정을 날씨에 비유한 내용으로 알맞지 \uline{\underline{않은}} 것은?}
\begin{kfChoices}
\kfChoice 흐린 날처럼 슬픈 날도 있다.
\kfChoice 감정은 날씨처럼 변할 수 있다.
\kfChoice 맑은 날씨처럼 항상 기뻐야만 한다.
\kfChoice 비가 오면 우산을 쓰듯 감정도 조절이 필요하다.
\end{kfChoices}
\end{kfQBlock}

\begin{kfQBlock}{04}{객관식}
\kfQStem{시 '내 마음의 색깔'에서 화가 나는 마음은 무슨 색이었나요?}
\begin{kfChoices}
\kfChoice 파랑
\kfChoice 노랑
\kfChoice 빨강
\kfChoice 초록
\end{kfChoices}
\end{kfQBlock}

\begin{kfQBlock}{05}{객관식}
\kfQStem{이 시에서 말하고자 하는 바와 가장 거리가 \uline{먼} 것은 무엇인가요?}
\begin{kfChoices}
\kfChoice 내 마음속에서 일어나는 여러 가지 감정들은 모두 가치가 있다.
\kfChoice 부정적인 감정이 들 때는 그것을 숨기거나 억지로 참아야 한다.
\kfChoice 오늘은 어떤 감정을 느끼더라도 그것을 있는 그대로 받아들이겠다.
\kfChoice 기쁨뿐만 아니라 슬픔이나 화나는 감정도 나의 자연스러운 일부이다.
\end{kfChoices}
\end{kfQBlock}

\begin{kfQBlock}{06}{객관식}
\kfQStem{화가 났을 때 마음을 지키는 방법으로 알맞지 \uline{\underline{않은}} 것은 무엇인가요?}
\begin{kfChoices}
\kfChoice 숫자를 셉니다.
\kfChoice 숨을 깊게 쉽니다.
\kfChoice 물건을 집어 던집니다.
\kfChoice 자신의 기분을 말합니다.
\end{kfChoices}
\end{kfQBlock}

\begin{kfQBlock}{07}{객관식}
\kfQStem{친구에게 내 마음을 가장 잘 전한 문장은 무엇인가요?}
\begin{kfChoices}
\kfChoice "너 정말 바보 같아."
\kfChoice "너 때문에 다 망쳤어."
\kfChoice "시끄러우니까 조용히 해."
\kfChoice "나는 네가 거짓말을 해서 화가 났어."
\end{kfChoices}
\end{kfQBlock}

\begin{kfQBlock}{08}{객관식}
\kfQStem{숨을 들이마시고 내쉴 때 우리 몸에 나타나는 변화는 무엇인가요?}
\begin{kfChoices}
\kfChoice 몸의 긴장이 풀립니다.
\kfChoice 갑자기 잠이 쏟아집니다.
\kfChoice 배가 풍선처럼 딱딱해집니다.
\kfChoice 목소리가 이전보다 더 커집니다.
\end{kfChoices}
\end{kfQBlock}

\begin{kfQBlock}{09}{객관식}
\kfQStem{다음 문장의 빈칸에 들어갈 알맞은 말은 무엇인가요?}
\begin{kfEvidence}발표를 잘 마치고 나니 마음이 \_\_\_\_\_\_.\end{kfEvidence}
\begin{kfChoices}
\kfChoice 화났다
\kfChoice 뿌듯했다
\kfChoice 무서웠다
\kfChoice 속상했다
\end{kfChoices}
\end{kfQBlock}

\begin{kfQBlock}{10}{서술형}
\kfQStem{다음 문장을 완성하세요.}
\begin{kfEvidence}"나는 친구와 싸웠을 때 마음이 \kfFillBlank[42mm]{}."\end{kfEvidence}
\kfAnswerNote{답안}{4}
\end{kfQBlock}




\clearpage

\kfChapterCover{3}{소쉬르2 (20주차) (3)}{Chapter 3}{이번 챕터의 학습 내용을 시작합니다.}
\begin{kfChapterAreaList}
\kfChapterAreaItem{어휘}{kfAreaVocab}{하루를 채우는 낱말 익히기}
\kfChapterAreaItem{개념}{kfAreaConcept}{왜 규칙적으로 생활해야 할까요?}
\kfChapterAreaItem{문학}{kfAreaLiterature}{「"민수야, 이제 잘 시간이야. 어서 자야…」}
\kfChapterAreaItem{비문학}{kfAreaNonlit}{건강한 하루를 만드는 습관}
\kfChapterAreaItem{실력확인}{kfAreaWeekly}{}
\end{kfChapterAreaList}

\kfStartVocab{하루를 채우는 낱말 익히기}


\begin{kfVocab}
\kfVocabItem{1}{앞으로 할 일을 미리 생각하여 정하는 것}
\kfVocabItem{2}{아침에 잠에서 깨어 일어나는 것}
\kfVocabItem{3}{오랫동안 되풀이하여 몸에 밴 행동}
\kfVocabItem{4}{날마다 규칙적으로 하는 일}
\kfVocabItem{5}{어지러운 물건을 제자리에 가지런히 치우는 것}
\kfVocabItem{6}{잠자리에 들어 잠을 자는 것}
\end{kfVocab}


\kfStartConcept{왜 규칙적으로 생활해야 할까요?}


\kfPassageNoteNT{%
하루를 건강하게 보내려면 규칙적인 생활이 필요해요.

제때 밥을 먹고 잠을 자면 몸이 튼튼해져요.

또, 잠을 푹 자고 일어나면 기분이 좋아져서 상쾌하게 하루를 시작할 수 있어요.

마지막으로 미리 계획을 세우면 허둥지둥하지 않고 중요한 일을 놓치지 않아요.
\par\medskip\begin{itemize}[leftmargin=18pt, itemsep=2pt, topsep=2pt, label=\textbullet]
\item 제때 식사 — 아침·점심·저녁 시간 지키기
\item 충분한 잠 — 매일 같은 시각에 자고 일어나기
\item 계획 세우기 — 학교 시간표·숙제 미리 정리
\end{itemize}
}

\kfPassageNote{일이 일어난 차례를 나타내는 말}{%
글을 쓰거나 말할 때 순서를 나타내는 말을 쓰면 내용을 더 잘 이해할 수 있어요. \\[14pt]
*  먼저(우선):가장 처음 할 일을 말할 때 \\[4pt]
*  그다음에(이어서):뒤이어 할 일을 말할 때 \\[4pt]
*  마지막으로:맨 끝에 할 일을 말할 때
}


\kfActivityBg \par\kfMaybeNeedspace{24\baselineskip}\noindent{\kfTipMark\,\,}{\small\color{kfInk} 위에서 배운 낱말 중 알맞은 것을 골라 빈칸을 채워 보세요.}\par\nopagebreak[4]\smallskip\nopagebreak[4]
\begin{kfBlankList}
\kfBlankItem 다 쓴 물건은 제자리에 \kfFillBlank[12mm]{}을 합니다.
\kfBlankItem 소풍 가서 보물찾기를 할 \kfFillBlank[12mm]{}을 세웠습니다.
\kfBlankItem 키가 크려면 밤 9시에는 \kfFillBlank[12mm]{}을 해야 합니다.
\kfBlankItem 아침에 일어나서 세수를 하는 것은 나의 \kfFillBlank[12mm]{}입니다.
\end{kfBlankList}

\kfActivityBg \par\kfMaybeNeedspace{24\baselineskip}\noindent{\kfTipMark\,\,}{\small\color{kfInk} 문제를 읽고 O, X 표시를 해 보세요.}\par\nopagebreak[4]\smallskip\nopagebreak[4]
\begin{kfOXList}
\kfOXItem 늦잠을 자는 것은 규칙적인 생활이다. \kfFillBlank[12mm]{} \kfOX
\kfOXItem 규칙적인 생활은 몸을 튼튼하게 해 준다. \kfFillBlank[12mm]{} \kfOX
\kfOXItem 계획을 세우면 중요한 일을 잊어버리기 쉽다. \kfFillBlank[12mm]{} \kfOX
\kfOXItem 잠을 푹 자야 기분 좋게 하루를 시작할 수 있다. \kfFillBlank[12mm]{} \kfOX
\kfOXItem 배가 고플 때만 밥을 먹는 것이 가장 좋은 습관이다. \kfFillBlank[12mm]{} \kfOX
\end{kfOXList}

\kfSecHeader{kfAreaConcept}{문제}{}

\begin{kfQBlock}{01}{객관식}
\kfQStem{다음 문장을 순서에 맞게 이어 주는 말로 연결해 보세요.}
\begin{kfChoices}
\kfChoice ㉠ 먼저 - ㉡ 마지막으로
\kfChoice ㉠ 먼저 - ㉡ 그다음에
\kfChoice ㉠ 마지막으로 - ㉡ 먼저
\kfChoice ㉠ 그다음에 - ㉡ 먼저
\end{kfChoices}
\end{kfQBlock}
\begin{kfQBlock}{02}{서술형}
\kfQStem{라면을 끓이는 순서를 생각하며 \kfFillBlank[12mm]{}에 알맞은 말을 넣으세요.}
\kfAnswerNote{답안}{4}
\end{kfQBlock}


\kfStartLit{「"민수야, 이제 잘 시간이야. 어서 자야…」}


\kfPassageNote{}{%
"민수야, 이제 잘 시간이야. 어서 자야지."

엄마가 방문을 열고 말씀하셨습니다. 하지만 민수는 침대에 눕기 싫었습니다.

"엄마, 10분만 더 놀다 잘게요. 로봇 조립이 거의 다 끝났단 말이에요."

민수는 시계를 보며 떼를 썼습니다. 시계바늘은 벌써 밤 10시를 넘어가고 있었습니다.

엄마는 한숨을 쉬며 나가셨고, 민수는 신나게 로봇을 가지고 놀았습니다.

'조금만 더, 조금만 더 해야지.'

어느새 시계는 12시를 가리키고 있었습니다. 민수는 그제야 하품을 크게 하고 잠이 들었습니다.

다음 날 아침이었습니다.

"민수야! 학교 늦겠다! 얼른 일어나!"

엄마의 목소리에 민수는 깜짝 놀라 일어났습니다. 머리가 띵하고 눈이 자꾸만 감겼습니다. 학교에 가서도 선생님 말씀이 귀에 들어오지 않았습니다. 꾸벅꾸벅 졸다가 짝꿍 지우가 쿡 찌르는 바람에 깜짝 놀라 깼습니다.

'아, 어제 일찍 잘걸. 너무 졸려.'

민수는 책상에 엎드리며 후회했습니다.
}


\kfActivityBg \par\kfMaybeNeedspace{24\baselineskip}\noindent{\kfTipMark\,\,}{\small\color{kfInk} 빈칸에 알맞은 말을 채워 보세요.}\par\nopagebreak[4]\smallskip\nopagebreak[4]
\begin{kfBlankList}
\kfBlankItem 이야기 속 사건을 일어난 차례대로 번호를 써 보세요.
\begin{kfEvidence}1) \kfFillBlank[12mm]{} 학교 수업 시간에 꾸벅꾸벅 졸았다. \\\relax
2) \kfFillBlank[12mm]{} 다음 날 아침, 늦잠을 자서 머리가 아팠다. \\\relax
3) \kfFillBlank[12mm]{} 엄마가 자라고 했지만 더 놀고 싶어서 늦게 잤다.\end{kfEvidence}
\end{kfBlankList}

\kfActivityBg \par\kfMaybeNeedspace{18\baselineskip}\noindent{\kfTipMark\,\,}{\small\color{kfInk} 주어진 안내에 맞춰 자유롭게 글로 표현해 보세요.}\par\nopagebreak[4]\smallskip\nopagebreak[4]
\kfWritingPrompt{잠을 자지 않아 힘들어하는 민수에게, 일찍 자면 어떤 점이 좋은지 알려 주는 짧은 편지를 써 보세요.}
\kfWriting{답안 작성}{8}

\kfSecHeader{kfAreaLiterature}{문제}{}

\begin{kfQBlock}{01}{객관식}
\kfQStem{민수가 잠자기 싫어했던 이유는 무엇인가요?}
\begin{kfChoices}
\kfChoice 저녁밥 먹은 배가 너무 고파서
\kfChoice 미뤄 둔 숙제를 다 못 해서
\kfChoice 보고 싶은 텔레비전을 못 봐서
\kfChoice 로봇 조립을 더 하고 싶어서
\end{kfChoices}
\end{kfQBlock}
\begin{kfQBlock}{02}{객관식}
\kfQStem{다음 날 아침 민수의 상태로 알맞은 것은?}
\begin{kfChoices}
\kfChoice 기분이 상쾌했다.
\kfChoice 밥맛이 아주 좋았다.
\kfChoice 머리가 띵하고 졸렸다.
\kfChoice 평소보다 일찍 일어났다.
\end{kfChoices}
\end{kfQBlock}
\begin{kfQBlock}{03}{객관식}
\kfQStem{민수가 학교에서 한 행동은 무엇인가요?}
\begin{kfChoices}
\kfChoice 꾸벅꾸벅 졸았다.
\kfChoice 발표를 씩씩하게 했다.
\kfChoice 친구들과 신나게 축구를 했다.
\kfChoice 선생님 말씀을 열심히 들었다.
\end{kfChoices}
\end{kfQBlock}
\begin{kfQBlock}{04}{객관식}
\kfQStem{민수가 마지막에 후회한 것은 무엇인가요?}
\begin{kfChoices}
\kfChoice 학교에 늦게 간 것
\kfChoice 아침밥을 먹지 않은 것
\kfChoice 어제 일찍 자지 않은 것
\kfChoice 로봇 조립을 완성하지 못한 것
\end{kfChoices}
\end{kfQBlock}
\begin{kfQBlock}{05}{객관식}
\kfQStem{민수가 건강하게 생활하려면 어떻게 해야 할까요?}
\begin{kfChoices}
\kfChoice 정해진 시간에 잠을 잔다.
\kfChoice 학교에서 낮잠을 많이 잔다.
\kfChoice 밤새도록 장난감을 가지고 논다.
\kfChoice 아침밥을 먹지 않고 학교에 간다.
\end{kfChoices}
\end{kfQBlock}


\kfStartNonlit{건강한 하루를 만드는 습관}


\kfPassageNote{건강한 하루를 만드는 습관}{%
우리가 매일 건강하게 지내려면 올바른 습관을 가지는 것이 매우 중요합니다. 아침부터 밤까지 우리가 실천할 수 있는 건강한 약속들을 순서대로 알아봅시다.

먼저 기분 좋은 아침을 시작해야 합니다. 아침에 눈을 뜨면 일찍 일어나서 몸을 쭉 펴는 기지개를 켭니다. 기지개는 잠들었던 우리 몸을 깨워 줍니다. 그리고 아침밥을 거르지 않고 꼭 먹어야 합니다. 아침밥을 든든하게 먹어야 학교에서 공부하고 뛰어놀 수 있는 힘이 생기기 때문입니다.

점심과 저녁 시간에도 지켜야 할 약속이 있습니다. 학교에서는 친구들과 즐겁게 공부하고 운동장에서 신나게 뛰어놉니다. 밥을 먹은 뒤에는 이를 깨끗이 닦습니다. 학교가 끝나고 집에 돌아오면 가장 먼저 손과 발을 깨끗이 씻습니다. 저녁밥을 먹은 후에는 가족들과 재미있게 이야기를 나눕니다.

마지막으로 밤에는 하루를 잘 마무리해야 합니다. 잠자리에 들기 전에 내일 학교에 가져갈 물건들을 미리 챙깁니다. 준비물을 미리 챙기면 마음 편히 쉴 수 있습니다. 그러고 나서 일찍 잠자리에 듭니다. 잠을 푹 자야 키도 쑥쑥 크고 다음 날 아침을 상쾌하게 맞이할 수 있습니다.
}


\kfActivityBg \par\needspace{15\baselineskip}
\kfSecHeader{kfAreaNonlit}{활동}{문장 독해}
\par\kfMaybeNeedspace{32\baselineskip}\noindent{\kfTipMark\,\,}{\small\color{kfInk} 각 문장을 읽고 물음에 답하세요.}\par\nopagebreak[4]\smallskip\nopagebreak[4]
\begin{kfSentenceUnit}{1}{몸을 우리 기지개는}
\kfSentQ{1.}{(1) 이 문장에서 '누가'에 해당하는 것은 무엇인가요?}
\begin{kfChoices}
\kfChoice 몸을
\kfChoice 우리
\kfChoice 기지개는
\end{kfChoices}
\kfSentQ{1.}{(2) 이 문장에서 '무엇을'에 해당하는 것은 무엇인가요?}
\begin{kfChoices}
\kfChoice 잠들었던
\kfChoice 기지개는
\kfChoice 우리 몸을
\end{kfChoices}
\kfSentQ{1.}{(3) 이 문장에서 '어찌하다'에 해당하는 것은 무엇인가요?}
\begin{kfChoices}
\kfChoice 깨운다
\kfChoice 기지개는
\kfChoice 우리 몸을
\end{kfChoices}
\end{kfSentenceUnit}

\begin{kfSentenceUnit}{2}{학교에서 점심밥을 우리는}
\kfSentQ{2.}{(1) 이 문장에서 '누가'에 해당하는 것은 무엇인가요?}
\begin{kfChoices}
\kfChoice 학교에서
\kfChoice 점심밥을
\kfChoice 우리는
\end{kfChoices}
\kfSentQ{2.}{(2) 이 문장에서 '무엇을'에 해당하는 것은 무엇인가요?}
\begin{kfChoices}
\kfChoice 먹는다.
\kfChoice 학교에서
\kfChoice 점심밥을
\end{kfChoices}
\kfSentQ{2.}{(3) 이 문장에서 '어찌하다'에 해당하는 것은 무엇인가요?}
\begin{kfChoices}
\kfChoice 점심밥을
\kfChoice 먹는다
\kfChoice 학교에서
\end{kfChoices}
\end{kfSentenceUnit}

\kfActivityBg \par\kfMaybeNeedspace{24\baselineskip}\noindent{\kfTipMark\,\,}{\small\color{kfInk} 빈칸에 알맞은 말을 채워 보세요.}\par\nopagebreak[4]\smallskip\nopagebreak[4]
\begin{kfBlankList}
\kfBlankItem 글을 읽고 빈칸에 알맞은 말을 넣어 내용을 정리해 봅시다.

매일 건강하게 지내려면 올바른 습관이 필요합니다.  아침에는 일찍 일어나서 기지개를 켜고 ①\kfFillBlank[12mm]{}을 꼭 먹어야 합니다.  낮에는 신나게 뛰어놀고 손과 발을 깨끗이 씻습니다.  밤에는 학교에 가져갈 ②\kfFillBlank[12mm]{}를 미리 챙긴 뒤 일찍 ③\kfFillBlank[12mm]{}에 들어야 합니다.
\end{kfBlankList}

\kfActivityBg \par\kfMaybeNeedspace{18\baselineskip}\noindent{\kfTipMark\,\,}{\small\color{kfInk} 주어진 안내에 맞춰 자유롭게 글로 표현해 보세요.}\par\nopagebreak[4]\smallskip\nopagebreak[4]
\kfWritingPrompt{글에서 배운 내용을 떠올리며 내가 꼭 지키고 싶은 건강 약속을 써 봅시다.}
\kfWriting{답안 작성}{8}

\kfSecHeader{kfAreaNonlit}{문제}{}

\begin{kfQBlock}{01}{객관식}
\kfQStem{이 글에서 아침에 일어나 기지개를 켜는 것의 효과로 알려 준 것은 무엇인가요?}
\begin{kfChoices}
\kfChoice 잠자리에 들기 전 준비물을 챙길 수 있다.
\kfChoice 학교에서 공부하고 뛰어놀 힘이 생긴다.
\kfChoice 잠들었던 우리 몸을 깨워 준다.
\kfChoice 다음 날 아침을 상쾌하게 맞이하게 된다.
\end{kfChoices}
\end{kfQBlock}
\begin{kfQBlock}{02}{객관식}
\kfQStem{이 글에서 아침밥을 거르지 않고 꼭 먹어야 하는 까닭으로 든 것은 무엇인가요?}
\begin{kfChoices}
\kfChoice 잠들었던 몸이 깨어나기 때문이다.
\kfChoice 학교에서 공부하고 뛰어놀 힘이 생기기 때문이다.
\kfChoice 키가 쑥쑥 자라기 때문이다.
\kfChoice 잠자리에서 마음 편히 쉴 수 있기 때문이다.
\end{kfChoices}
\end{kfQBlock}
\begin{kfQBlock}{03}{객관식}
\kfQStem{이 글에서 학교가 끝나고 집에 돌아오면 '가장 먼저' 하라고 한 일은 무엇인가요?}
\begin{kfChoices}
\kfChoice 손과 발을 깨끗이 씻는다.
\kfChoice 가족들과 재미있게 이야기를 나눈다.
\kfChoice 잠자리에 들기 전에 준비물을 챙긴다.
\kfChoice 운동장에서 신나게 뛰어논다.
\end{kfChoices}
\end{kfQBlock}
\begin{kfQBlock}{04}{객관식}
\kfQStem{이 글에서 잠자리에 들기 전 준비물을 미리 챙겨야 하는 까닭으로 든 것은 무엇인가요?}
\begin{kfChoices}
\kfChoice 마음 편히 쉴 수 있기 때문이다.
\kfChoice 잠들었던 몸을 깨울 수 있기 때문이다.
\kfChoice 아침밥을 더 많이 먹기 위해서이다.
\kfChoice 손과 발을 깨끗이 씻기 위해서이다.
\end{kfChoices}
\end{kfQBlock}
\begin{kfQBlock}{05}{객관식}
\kfQStem{이 글에서 밤에 잠을 푹 자야 한다고 한 까닭으로 든 것이 \uline{\underline{아닌}} 것은 무엇인가요?}
\begin{kfChoices}
\kfChoice 키도 쑥쑥 크기 때문이다.
\kfChoice 다음 날 학교에 지각하지 않기 때문이다.
\kfChoice 다음 날 아침을 상쾌하게 맞이할 수 있기 때문이다.
\kfChoice 일찍 일어나 기지개를 켜야 하기 때문이다.
\end{kfChoices}
\end{kfQBlock}


\kfStartWeekly{실력확인 영역 학습}


\noindent\textit{\small 제한시간: 20분 / 총 10문항}\par\medskip
\par\noindent\textbf{[4\textasciitilde{}5] 다음 글을 읽고 물음에 답하시오.}\par\medskip
\kfPassageNote{지문 4}{%
"민수야, 이제 잘 시간이야. 어서 자야지." 엄마가 방문을 열고 말씀하셨습니다. 하지만 민수는 침대에 눕기 싫었습니다. "엄마, 10분만 더 놀다 잘게요. 로봇 조립이 거의 다 끝났단 말이에요." 민수는 시계를 보며 떼를 썼습니다. 시계바늘은 벌써 밤 10시를 넘어가고 있었습니다. 엄마는 한숨을 쉬며 나가셨고, 민수는 신나게 로봇을 가지고 놀았습니다. '조금만 더, 조금만 더 해야지.' 어느새 시계는 12시를 가리키고 있었습니다. 민수는 그제야 하품을 크게 하고 잠이 들었습니다. 다음 날 아침이었습니다. "민수야! 학교 늦겠다! 얼른 일어나!" 엄마의 목소리에 민수는 깜짝 놀라 일어났습니다. 머리가 띵하고 눈이 자꾸만 감겼습니다. 학교에 가서도 선생님 말씀이 귀에 들어오지 않았습니다. 꾸벅꾸벅 졸다가 짝꿍 지우가 쿡 찌르는 바람에 깜짝 놀라 깼습니다. '아, 어제 일찍 잘걸. 너무 졸려.' 민수는 책상에 엎드리며 후회했습니다.
}
\par\noindent\textbf{[6\textasciitilde{}7] 다음 글을 읽고 물음에 답하시오.}\par\medskip
\kfPassageNote{지문 6}{%
우리가 매일 건강하게 지내려면 올바른 습관을 가지는 것이 매우 중요합니다. 아침부터 밤까지 우리가 실천할 수 있는 건강한 약속들을 순서대로 알아봅시다.

먼저 기분 좋은 아침을 시작해야 합니다. 아침에 눈을 뜨면 일찍 일어나서 몸을 쭉 펴는 기지개를 켭니다. 기지개는 잠들었던 우리 몸을 깨워 줍니다. 그리고 아침밥을 거르지 않고 꼭 먹어야 합니다. 아침밥을 든든하게 먹어야 학교에서 공부하고 뛰어놀 수 있는 힘이 생기기 때문입니다.

점심과 저녁 시간에도 지켜야 할 약속이 있습니다. 학교에서는 친구들과 즐겁게 공부하고 운동장에서 신나게 뛰어놉니다. 밥을 먹은 뒤에는 이를 깨끗이 닦습니다. 학교가 끝나고 집에 돌아오면 가장 먼저 손과 발을 깨끗이 씻습니다. 저녁밥을 먹은 후에는 가족들과 재미있게 이야기를 나눕니다.

마지막으로 밤에는 하루를 잘 마무리해야 합니다. 잠자리에 들기 전에 내일 학교에 가져갈 물건들을 미리 챙깁니다. 준비물을 미리 챙기면 마음 편히 쉴 수 있습니다. 그러고 나서 일찍 잠자리에 듭니다. 잠을 푹 자야 키도 쑥쑥 크고 다음 날 아침을 상쾌하게 맞이할 수 있습니다.
}
\begin{kfQBlock}{01}{객관식}
\kfQStem{날마다 규칙적으로 하는 일을 뜻하는 낱말은?}
\begin{kfChoices}
\kfChoice 일과
\kfChoice 소풍
\kfChoice 차례
\kfChoice 늦잠
\end{kfChoices}
\end{kfQBlock}

\begin{kfQBlock}{02}{객관식}
\kfQStem{'습관'과 뜻이 비슷한 낱말은?}
\begin{kfChoices}
\kfChoice 버릇
\kfChoice 기적
\kfChoice 선물
\kfChoice 장난
\end{kfChoices}
\end{kfQBlock}

\begin{kfQBlock}{03}{객관식}
\kfQStem{계획을 세우면 좋은 점은 무엇인가요?}
\begin{kfChoices}
\kfChoice 늦잠을 잘 수 있다.
\kfChoice 할 일을 미루게 된다.
\kfChoice 친구와 더 많이 놀 수 있다.
\kfChoice 중요한 일을 잊지 않고 챙길 수 있다.
\end{kfChoices}
\end{kfQBlock}

\begin{kfQBlock}{04}{객관식}
\kfQStem{'잠자기 싫은 민수' 이야기에서 민수가 늦잠을 잔 까닭은?}
\begin{kfChoices}
\kfChoice 몸이 아파서
\kfChoice 숙제가 많아서
\kfChoice 늦게까지 놀아서
\kfChoice 알람 시계가 고장 나서
\end{kfChoices}
\end{kfQBlock}

\begin{kfQBlock}{05}{객관식}
\kfQStem{다음 날 아침, 학교에 간 민수의 모습으로 알맞은 것은 무엇인가요?}
\begin{kfChoices}
\kfChoice 기분이 상쾌했다.
\kfChoice 밥맛이 아주 좋았다.
\kfChoice 친구들과 신나게 놀았다.
\kfChoice 머리가 띵하고 자꾸 졸렸다.
\end{kfChoices}
\end{kfQBlock}

\begin{kfQBlock}{06}{객관식}
\kfQStem{다음 중 하루 일과에서 가장 먼저 일어난 일은?}
\begin{kfChoices}
\kfChoice 기지개를 켰다.
\kfChoice 이를 깨끗이 닦았다.
\kfChoice 운동장에서 뛰어놀았다.
\kfChoice 학교에 가져갈 준비물을 챙겼다.
\end{kfChoices}
\end{kfQBlock}

\begin{kfQBlock}{07}{객관식}
\kfQStem{건강한 생활 습관이 \uline{\underline{아닌}} 것은?}
\begin{kfChoices}
\kfChoice 기지개 켜기
\kfChoice 손과 발 씻기
\kfChoice 늦게 잠자리에 들기
\kfChoice 내일 학교 준비물 챙기기
\end{kfChoices}
\end{kfQBlock}

\begin{kfQBlock}{08}{객관식}
\kfQStem{일이 일어나는 순서를 나타낼 때 쓰는 말이 \uline{\underline{아닌}} 것은?}
\begin{kfChoices}
\kfChoice 먼저
\kfChoice 예쁘게
\kfChoice 그다음에
\kfChoice 마지막으로
\end{kfChoices}
\end{kfQBlock}

\begin{kfQBlock}{09}{객관식}
\kfQStem{다음 문장의 빈칸에 들어갈 알맞은 말은?}
\begin{kfEvidence}"씨앗을 심었습니다. \kfFillBlank[42mm]{}, 물을 주었습니다."\end{kfEvidence}
\begin{kfChoices}
\kfChoice 먼저
\kfChoice 하지만
\kfChoice 그다음에
\kfChoice 왜냐하면
\end{kfChoices}
\end{kfQBlock}

\begin{kfQBlock}{10}{객관식}
\kfQStem{다음 문장을 읽고 주어, 목적어, 서술어를 바르게 짝지은 것은 무엇인가요?}
\begin{kfEvidence}기지개는 잠들었던 우리 몸을 깨운다.\end{kfEvidence}
\begin{kfChoices}
\kfChoice 기지개는 / 잠들었던 / 깨운다
\kfChoice 잠들었던 / 우리 몸을 / 깨운다
\kfChoice 기지개는 / 우리 몸을 / 깨운다
\kfChoice 기지개는 / 잠들었던 우리 몸을 / 깨운다
\end{kfChoices}
\end{kfQBlock}




\clearpage

\kfChapterCover{4}{소쉬르2 (20주차) (4)}{Chapter 4}{이번 챕터의 학습 내용을 시작합니다.}
\begin{kfChapterAreaList}
\kfChapterAreaItem{어휘}{kfAreaVocab}{동네를 채우는 낱말 익히기}
\kfChapterAreaItem{개념}{kfAreaConcept}{내용을 간추린다는 건 뭘까요?}
\kfChapterAreaItem{문학}{kfAreaLiterature}{「어느 화창한 토요일, 민지와 수호는 동네…」}
\kfChapterAreaItem{비문학}{kfAreaNonlit}{도서관에서 지켜야 할 약속}
\kfChapterAreaItem{실력확인}{kfAreaWeekly}{}
\end{kfChapterAreaList}

\kfStartVocab{동네를 채우는 낱말 익히기}


\begin{kfVocab}
\kfVocabItem{1}{여러 사람이 함께 이용하는 곳 (예: 공원, 도서관, 지하철)}
\kfVocabItem{2}{사람들이 모여 사는 곳}
\kfVocabItem{3}{경치가 아름답거나 유명 해서 사람들이 많이 찾는 곳}
\kfVocabItem{4}{도서관이나 놀이터처럼 여러 사람이 이용하도록 만든 건물이나 곳}
\kfVocabItem{5}{가까이에 사는 집이나 사람들}
\kfVocabItem{6}{땅의 모양이나 길, 건물의 위치를 줄여서 종이에 그린 그림}
\end{kfVocab}


\kfStartConcept{내용을 간추린다는 건 뭘까요?}


\kfPassageNoteNT{%
긴 글이나 이야기를 읽고 '간추린다'는 것은 무슨 뜻일까요?

먼저, 긴 이야기를 짧게 줄여서 말하는 것을 뜻해요.

간추릴 때는 필요 없는 말은 빼고 가장 중요하고 알맹이가 되는 내용만 남겨요.

이렇게 내용을 간추리면 긴 글도 쉽게 기억할 수 있고, 다른 사람에게 핵심을 잘 전달할 수 있어요.
\par\medskip\begin{itemize}[leftmargin=18pt, itemsep=2pt, topsep=2pt, label=\textbullet]
\item 줄이기 — 긴 동화 → 한 문장 요약
\item 남기기 — 핵심 사건과 등장인물만
\item 빼기 — 반복되는 말, 꾸미는 말 생략
\end{itemize}
}

\kfPassageNote{내용을 간추리는 방법}{%
글을 읽고 중요한 내용을 찾아서 간단하게 줄이는 연습을 해 봐요. \\[14pt]
* 누가: 이야기의 주인공 찾기 \\[4pt]
* 어디서: 일이 일어난 장소 찾기 \\[4pt]
* 무엇을: 어떤 일이 일어났는지 찾기 \\[4pt]
* 그래서: 결과가 어떻게 되었는지 찾기
}


\kfActivityBg \par\kfMaybeNeedspace{24\baselineskip}\noindent{\kfTipMark\,\,}{\small\color{kfInk} 위에서 배운 낱말 중 알맞은 것을 골라 빈칸을 채워 보세요.}\par\nopagebreak[4]\smallskip\nopagebreak[4]
\begin{kfBlankList}
\kfBlankItem 나는 정이 넘치는 우리 \kfFillBlank[12mm]{}를 참 좋아합니다.
\kfBlankItem 길을 모를 때는 \kfFillBlank[12mm]{}를 보면 쉽게 찾을 수 있습니다.
\kfBlankItem 도서관 같은 \kfFillBlank[12mm]{}에서는 큰 소리로 떠들면 안 됩니다.
\kfBlankItem 옆집에 사는 \kfFillBlank[12mm]{}을/를 만나면 반갑게 인사합니다.
\end{kfBlankList}

\kfActivityBg \par\kfMaybeNeedspace{24\baselineskip}\noindent{\kfTipMark\,\,}{\small\color{kfInk} 문제를 읽고 O, X 표시를 해 보세요.}\par\nopagebreak[4]\smallskip\nopagebreak[4]
\begin{kfOXList}
\kfOXItem 도서관은 책을 읽거나 빌리는 곳이다. \kfFillBlank[12mm]{} \kfOX
\kfOXItem 지도를 보면 건물의 위치를 알 수 있다. \kfFillBlank[12mm]{} \kfOX
\kfOXItem 내용을 간추릴 때는 모든 문장을 다 써야 한다. \kfFillBlank[12mm]{} \kfOX
\kfOXItem 공공장소에서는 나 혼자 마음대로 행동해도 된다. \kfFillBlank[12mm]{} \kfOX
\kfOXItem 간추리기는 중요한 내용만 남기고 짧게 줄이는 것이다. \kfFillBlank[12mm]{} \kfOX
\end{kfOXList}

\kfActivityBg \par\kfMaybeNeedspace{24\baselineskip}\noindent{\kfTipMark\,\,}{\small\color{kfInk} 다음 활동을 풀어 보세요.}\par\nopagebreak[4]\smallskip\nopagebreak[4]
\begin{kfBlankList}
\kfBlankItem 다음 긴 문장을 짧게 간추린 것으로 알맞은 것은?
\begin{kfEvidence}"나는 어제 학교 운동장에서 친구들과 함께 땀을 흘리며 신나게 축구를 했습니다."

→ (간추린 문장): 나는 어제 학교에서 \kfFillBlank[12mm]{}.\end{kfEvidence}
\end{kfBlankList}

\kfSecHeader{kfAreaConcept}{문제}{}

\begin{kfQBlock}{01}{객관식}
\kfQStem{다음 문장에서 가장 중요한 낱말(핵심어)을 찾아보세요.}
\begin{kfChoices}
\kfChoice 우리
\kfChoice 동네
\kfChoice 도서관
\kfChoice 가득합니다
\end{kfChoices}
\end{kfQBlock}


\kfStartLit{「어느 화창한 토요일, 민지와 수호는 동네…」}


\kfPassageNote{}{%
어느 화창한 토요일, 민지와 수호는 동네 놀이터에서 낡은 종이 한 장을 주웠습니다.

"어? 이게 뭐지? 보물 지도인가 봐!"

종이에는 우리 동네 그림과 함께 '보물이 있는 곳'이라는 표시가 되어 있었습니다. 민지와 수호는 두근거리는 마음으로 지도를 따라갔습니다.

가장 먼저 '행복 빵집'이 나왔습니다. 고소한 빵 냄새가 났지만, 보물은 없었습니다.

"지도를 보니까 우체국을 지나서 오른쪽으로 가야 해."

수호가 손가락으로 길을 가리켰습니다. 둘은 우체국을 지나 골목길로 들어갔습니다. 담벼락에는 예쁜 꽃 그림이 그려져 있었습니다.

"와, 우리 동네에 이런 예쁜 길이 있었네?"

민지가 감탄하며 말했습니다.

드디어 지도에 표시된 장소인 '초록 공원'에 도착했습니다. 커다란 느티나무 아래에 작은 상자가 놓여 있었습니다. 상자 안에는 거울 하나와 쪽지가 들어 있었습니다. 쪽지에는 이렇게 적혀 있었습니다.

<진짜 보물은 바로 우리 동네를 사랑하는 너희들이야!>

민지와 수호는 서로의 얼굴을 보며 활짝 웃었습니다.
}


\kfActivityBg \par\kfMaybeNeedspace{24\baselineskip}\noindent{\kfTipMark\,\,}{\small\color{kfInk} 빈칸에 알맞은 말을 채워 보세요.}\par\nopagebreak[4]\smallskip\nopagebreak[4]
\begin{kfBlankList}
\kfBlankItem 이 이야기를 '일이 일어난 차례'대로 간추려 보세요. 빈칸에 들어갈 말은?
\begin{kfEvidence}(1) 놀이터에서 보물 지도를 주웠다. \\\relax
(2) \kfFillBlank[12mm]{}을/를 지나 골목길을 걸었다. \\\relax
(3) 공원 느티나무 아래에서 상자를 발견했다.\end{kfEvidence}
\end{kfBlankList}

\kfActivityBg \par\kfMaybeNeedspace{18\baselineskip}\noindent{\kfTipMark\,\,}{\small\color{kfInk} 주어진 안내에 맞춰 자유롭게 글로 표현해 보세요.}\par\nopagebreak[4]\smallskip\nopagebreak[4]
\kfWritingPrompt{우리 동네에서 내가 가장 좋아하는 '보물 장소'는 어디인가요? 그곳을 좋아하는 이유와 함께 짧은 글로 소개해 보세요.}
\kfWriting{답안 작성}{8}

\kfSecHeader{kfAreaLiterature}{문제}{}

\begin{kfQBlock}{01}{객관식}
\kfQStem{민지와 수호가 놀이터에서 주운 것은 무엇인가요?}
\begin{kfChoices}
\kfChoice 예쁜 꽃
\kfChoice 낡은 종이
\kfChoice 맛있는 빵
\kfChoice 장난감 로봇
\end{kfChoices}
\end{kfQBlock}
\begin{kfQBlock}{02}{객관식}
\kfQStem{아이들이 가장 먼저 지나간 곳은 어디인가요?}
\begin{kfChoices}
\kfChoice 병원
\kfChoice 경찰서
\kfChoice 소방서
\kfChoice 행복 빵집
\end{kfChoices}
\end{kfQBlock}
\begin{kfQBlock}{03}{객관식}
\kfQStem{상자 안에 들어 있던 물건은 무엇인가요?}
\begin{kfChoices}
\kfChoice 지도
\kfChoice 사탕
\kfChoice 금동전
\kfChoice 거울과 쪽지
\end{kfChoices}
\end{kfQBlock}
\begin{kfQBlock}{04}{객관식}
\kfQStem{쪽지에서 말한 '진짜 보물'은 누구인가요?}
\begin{kfChoices}
\kfChoice 공원에 서 있는 큰 느티나무
\kfChoice 동네를 살피시는 이장님
\kfChoice 빵을 굽는 빵집 주인아저씨
\kfChoice 동네를 사랑하는 우리 아이들
\end{kfChoices}
\end{kfQBlock}
\begin{kfQBlock}{05}{객관식}
\kfQStem{민지와 수호가 보물을 찾으러 다니며 느낀 점으로 알맞은 것은?}
\begin{kfChoices}
\kfChoice 빵을 먹지 못해 배가 고팠다.
\kfChoice 길이 너무 복잡해서 무서웠다.
\kfChoice 다리가 아파서 짜증이 났다.
\kfChoice 새로운 모습을 봐서 즐거웠다.
\end{kfChoices}
\end{kfQBlock}


\kfStartNonlit{도서관에서 지켜야 할 약속}


\kfPassageNote{도서관에서 지켜야 할 약속}{%
도서관은 누구나 자유롭게 이용하는 공공 장소입니다. 도서관은 수많은 책이 모여 있는 지식의 보물 창고와 같습니다. 도서관에서 책을 읽을 때는 지켜야 할 약속이 있습니다. 첫째, 도서관 안에서는 조용히 해야 합니다. 도서관은 많은 사람이 모여 집중해서 책을 읽는 곳입니다. 다른 사람의 독서를 방해하지 않도록 발소리를 죽여 살금살금 걷고, 꼭 필요한 말이 있다면 아주 작은 목소리로 속삭여야 합니다. 둘째, 책을 소중히 다루어야 합니다. 도서관 책은 여러 사람이 함께 보는 물건입니다. 책을 찢거나 낙서를 하지 않고 깨끗하게 읽어야 합니다. 책을 다 보고 나서 아무 곳에나 두면 다른 사람이 그 책을 찾을 수 없게 됩니다. 다 읽은 책은 도서 운반 수레에 가지런히 올려두어 잘 정리되게 돕습니다.
}


\kfActivityBg \par\needspace{15\baselineskip}
\kfSecHeader{kfAreaNonlit}{활동}{문장 독해}
\par\kfMaybeNeedspace{32\baselineskip}\noindent{\kfTipMark\,\,}{\small\color{kfInk} 각 문장을 읽고 물음에 답하세요.}\par\nopagebreak[4]\smallskip\nopagebreak[4]
\begin{kfSentenceUnit}{1}{도서관을 이용한다. 사람들이}
\kfSentQ{1.}{(1) 이 문장에서 '누가'에 해당하는 것은 무엇인가요?}
\begin{kfChoices}
\kfChoice 도서관을
\kfChoice 이용한다.
\kfChoice 사람들이
\end{kfChoices}
\kfSentQ{1.}{(2) 이 문장에서 '무엇을'에 해당하는 것은 무엇인가요?}
\begin{kfChoices}
\kfChoice 사람들이
\kfChoice 이용한다.
\kfChoice 도서관을
\end{kfChoices}
\kfSentQ{1.}{(3) 이 문장에서 '어찌하다'에 해당하는 것은 무엇인가요?}
\begin{kfChoices}
\kfChoice 도서관을
\kfChoice 이용한다
\kfChoice 사람들이
\end{kfChoices}
\end{kfSentenceUnit}

\begin{kfSentenceUnit}{2}{지켜야 학생이 장소에서}
\kfSentQ{2.}{(1) 이 문장에서 '누가'에 해당하는 것은 무엇인가요?}
\begin{kfChoices}
\kfChoice 지켜야
\kfChoice 학생이
\kfChoice 장소에서
\end{kfChoices}
\kfSentQ{2.}{(2) 이 문장에서 '무엇을'에 해당하는 것은 무엇인가요?}
\begin{kfChoices}
\kfChoice 공공 장소에서
\kfChoice 약속을
\kfChoice 지켜야
\end{kfChoices}
\kfSentQ{2.}{(3) 이 문장에서 '어찌하다'에 해당하는 것은 무엇인가요?}
\begin{kfChoices}
\kfChoice 장소에서
\kfChoice 약속을
\kfChoice 배운다
\end{kfChoices}
\end{kfSentenceUnit}

\kfActivityBg \par\kfMaybeNeedspace{24\baselineskip}\noindent{\kfTipMark\,\,}{\small\color{kfInk} 빈칸에 알맞은 말을 채워 보세요.}\par\nopagebreak[4]\smallskip\nopagebreak[4]
\begin{kfBlankList}
\kfBlankItem 글을 읽고 빈칸에 알맞은 말을 넣어 문장을 완성해 봅시다.

도서관은 누구나 이용하는 ( ① ) 장소입니다. 이곳에서는 다른 사람을 방해하지 않도록 ( ② ) 해야 합니다. 또한 여러 사람이 함께 보는 책을 소중히 다루어야 합니다. 다 읽은 책은 다른 사람이 잘 찾을 수 있게 ( ③ ) 수레에 올려둡니다.
\end{kfBlankList}

\kfActivityBg \par\kfMaybeNeedspace{18\baselineskip}\noindent{\kfTipMark\,\,}{\small\color{kfInk} 주어진 안내에 맞춰 자유롭게 글로 표현해 보세요.}\par\nopagebreak[4]\smallskip\nopagebreak[4]
\kfWritingPrompt{내가 도서관을 이용할 때 가장 잘 지킬 수 있는 약속 한 가지를 정해서 써 봅시다.}
\kfWriting{답안 작성}{8}

\kfSecHeader{kfAreaNonlit}{문제}{}

\begin{kfQBlock}{01}{객관식}
\kfQStem{이 글에서 도서관을 어떻게 설명했나요?}
\begin{kfChoices}
\kfChoice 한 사람이 빌려가면 다른 사람은 들어갈 수 없는 곳
\kfChoice 학교 학생들만 이용할 수 있는 비밀 보관실
\kfChoice 책을 사고팔 수 있는 도서 가게
\kfChoice 누구나 자유롭게 이용하는 공공 장소
\end{kfChoices}
\end{kfQBlock}
\begin{kfQBlock}{02}{객관식}
\kfQStem{이 글에서 도서관 안에서 걸을 때 어떻게 하라고 했나요?}
\begin{kfChoices}
\kfChoice 책장 사이를 빠르게 뛰어다닌다.
\kfChoice 친구와 함께 박자를 맞춰 걷는다.
\kfChoice 발소리를 죽여 살금살금 걷는다.
\kfChoice 다 읽은 책을 끌면서 걷는다.
\end{kfChoices}
\end{kfQBlock}
\begin{kfQBlock}{03}{객관식}
\kfQStem{이 글에서 도서관 안에서 꼭 필요한 말을 해야 할 때 어떻게 하라고 했나요?}
\begin{kfChoices}
\kfChoice 책장 옆에 서서 책을 소리 내어 읽는다.
\kfChoice 다른 사람의 자리를 찾아가 또박또박 말한다.
\kfChoice 아주 작은 목소리로 속삭인다.
\kfChoice 도서 운반 수레 옆에서 친구와 큰 소리로 의논한다.
\end{kfChoices}
\end{kfQBlock}
\begin{kfQBlock}{04}{객관식}
\kfQStem{이 글에서 도서관 책을 소중히 다루는 방법으로 알려 준 것은 무엇인가요?}
\begin{kfChoices}
\kfChoice 책의 표지에 이름을 적어 둔다.
\kfChoice 자주 보는 쪽을 접어 표시해 둔다.
\kfChoice 책을 찢거나 낙서를 하지 않고 깨끗하게 읽는다.
\kfChoice 다 본 책은 살금살금 걸어 책장 안쪽에 끼워 둔다.
\end{kfChoices}
\end{kfQBlock}
\begin{kfQBlock}{05}{객관식}
\kfQStem{다 읽은 책은 어디에 두어야 하나요?}
\begin{kfChoices}
\kfChoice 책상 위
\kfChoice 바닥 앞
\kfChoice 창틀 옆
\kfChoice 도서 운반 수레
\end{kfChoices}
\end{kfQBlock}


\kfStartWeekly{실력확인 영역 학습}


\noindent\textit{\small 제한시간: 20분 / 총 10문항}\par\medskip
\par\noindent\textbf{[4\textasciitilde{}5] 다음 글을 읽고, 물음에 답하세요.}\par\medskip
\kfPassageNote{지문 4}{%
어느 화창한 토요일, 민지와 수호는 동네 놀이터에서 낡은 종이 한 장을 주웠습니다. "어? 이게 뭐지? 보물 지도인가 봐!" 종이에는 우리 동네 그림과 함께 '보물이 있는 곳'이라는 표시가 되어 있었습니다. 민지와 수호는 두근거리는 마음으로 지도를 따라갔습니다. 가장 먼저 '행복 빵집'이 나왔습니다. 고소한 빵 냄새가 났지만, 보물은 없었습니다. "지도를 보니까 우체국을 지나서 오른쪽으로 가야 해." 수호가 손가락으로 길을 가리켰습니다. 둘은 우체국을 지나 골목길로 들어갔습니다. 담벼락에는 예쁜 꽃 그림이 그려져 있었습니다. "와, 우리 동네에 이런 예쁜 길이 있었네?" 민지가 감탄하며 말했습니다. 드디어 지도에 표시된 장소인 '초록 공원'에 도착했습니다. 커다란 느티나무 아래에 작은 상자가 놓여 있었습니다. 상자 안에는 거울 하나와 쪽지가 들어 있었습니다. 쪽지에는 이렇게 적혀 있었습니다. <진짜 보물은 바로 우리 동네를 사랑하는 너희들이야!> 민지와 수호는 서로의 얼굴을 보며 활짝 웃었습니다.
}
\par\noindent\textbf{[6\textasciitilde{}7] 다음 글을 읽고, 물음에 답하세요.}\par\medskip
\kfPassageNote{지문 6}{%
도서관은 누구나 자유롭게 이용하는 공공 장소입니다. 도서관은 수많은 책이 모여 있는 지식의 보물 창고와 같습니다. 도서관에서 책을 읽을 때는 지켜야 할 약속이 있습니다. 첫째, 도서관 안에서는 조용히 해야 합니다. 도서관은 많은 사람이 모여 집중해서 책을 읽는 곳입니다. 다른 사람의 독서를 방해하지 않도록 발소리를 죽여 살금살금 걷고, 꼭 필요한 말이 있다면 아주 작은 목소리로 속삭여야 합니다. 둘째, 책을 소중히 다루어야 합니다. 도서관 책은 여러 사람이 함께 보는 물건입니다. 책을 찢거나 낙서를 하지 않고 깨끗하게 읽어야 합니다. 책을 다 보고 나서 아무 곳에나 두면 다른 사람이 그 책을 찾을 수 없게 됩니다. 다 읽은 책은 도서 운반 수레에 가지런히 올려두어 잘 정리되게 돕습니다.
}
\begin{kfQBlock}{01}{객관식}
\kfQStem{'매우 드물고 귀중한 물건'을 뜻하는 낱말은 무엇인가요?}
\begin{kfChoices}
\kfChoice 지도
\kfChoice 보물
\kfChoice 동네
\kfChoice 쓰레기
\end{kfChoices}
\end{kfQBlock}

\begin{kfQBlock}{02}{객관식}
\kfQStem{우리 동네를 소개하는 글로 알맞지 \uline{\underline{않은}} 것은 무엇인가요?}
\begin{kfChoices}
\kfChoice 나는 우리 동네가 싫습니다.
\kfChoice 우리 동네 도서관은 책이 많습니다.
\kfChoice 우리 동네에는 큰 공원이 있습니다.
\kfChoice 우리 동네 시장은 시끌벅적합니다.
\end{kfChoices}
\end{kfQBlock}

\begin{kfQBlock}{03}{객관식}
\kfQStem{다음 중 '공공장소'가 \uline{\underline{아닌}} 곳은?}
\begin{kfChoices}
\kfChoice 공원
\kfChoice 거실
\kfChoice 지하철
\kfChoice 도서관
\end{kfChoices}
\end{kfQBlock}

\begin{kfQBlock}{04}{객관식}
\kfQStem{이야기 '보물찾기 대작전'에서 보물이 있던 장소는 어디인가요?}
\begin{kfChoices}
\kfChoice 우체국
\kfChoice 행복 빵집
\kfChoice 초록 공원
\kfChoice 학교 운동장
\end{kfChoices}
\end{kfQBlock}

\begin{kfQBlock}{05}{객관식}
\kfQStem{상자 안에 금은보화 대신 '거울'이 들어 있었던 까닭은 무엇일까요?}
\begin{kfChoices}
\kfChoice 선물이어서
\kfChoice 보물이 없어서
\kfChoice 얼굴을 닦으라고
\kfChoice 진짜 보물인 자신의 모습을 보라고
\end{kfChoices}
\end{kfQBlock}

\begin{kfQBlock}{06}{객관식}
\kfQStem{도서관에서 책을 찢으면 안 되는 이유는 무엇인가요?}
\begin{kfChoices}
\kfChoice 종이가 아까워서
\kfChoice 함께 보는 책이니까
\kfChoice 내가 가질 책이 아니라서
\kfChoice 사서 선생님께 혼나니까
\end{kfChoices}
\end{kfQBlock}

\begin{kfQBlock}{07}{객관식}
\kfQStem{도서관에서 지켜야 할 약속으로 알맞은 것은 무엇인가요?}
\begin{kfChoices}
\kfChoice 뛰기
\kfChoice 장난치기
\kfChoice 조용히 하기
\kfChoice 큰 소리로 말하기
\end{kfChoices}
\end{kfQBlock}

\begin{kfQBlock}{08}{객관식}
\kfQStem{내용을 간추리는 방법으로 알맞지 \uline{\underline{않은}} 것은?}
\begin{kfChoices}
\kfChoice 중요한 내용을 찾는다.
\kfChoice 필요 없는 말은 뺀다.
\kfChoice 간단하게 줄여서 쓴다.
\kfChoice 책에 있는 모든 글자를 그대로 쓴다.
\end{kfChoices}
\end{kfQBlock}

\begin{kfQBlock}{09}{객관식}
\kfQStem{"토끼가 숲속에서 맛있는 풀을 먹었습니다."에서 일이 일어난 장소는 어디인가요?}
\begin{kfChoices}
\kfChoice 풀
\kfChoice 숲속
\kfChoice 운동장
\kfChoice 할머니 댁
\end{kfChoices}
\end{kfQBlock}

\begin{kfQBlock}{10}{객관식}
\kfQStem{다음 문장에서 꾸며 주는 말이나 구체적인 상황 설명을 빼고, ‘누가 무엇을 했다’는 중심 내용만 가장 잘 간추린 것은 무엇인가요?}
\begin{kfEvidence}"어제 우리 반 친구들은 학교 뒤뜰에 있는 텃밭에서 빨갛게 잘 익은 방울토마토를 기쁜 마음으로 땄습니다."\end{kfEvidence}
\begin{kfChoices}
\kfChoice 텃밭이 학교 뒤에 있다.
\kfChoice 친구들의 마음이 기뻤다.
\kfChoice 방울토마토가 빨갛게 익었다.
\kfChoice 친구들이 방울토마토를 땄다.
\end{kfChoices}
\end{kfQBlock}




\end{document}
